\chapter{Probabilistic and Objective-Shaped Memory: Kalman Attention, Gated KalmaNet, Next-Latent Prediction}\label{ch:probabilistic}
\chaptermark{Probabilistic and Objective-Shaped Memory}% short running header

Every recurrent memory in the previous chapters answers the same question. A stream of
key--value pairs $(\vk_1,\vv_1),(\vk_2,\vv_2),\dots$ arrives; which single matrix
$\mS_t \in \R^{d_v \times d_k}$ should the layer keep so that $\mS_t\vq$ answers later queries?
Linear attention (\cref{ch:linear}) adds outer products. The delta rule (\cref{ch:delta}) takes
one gradient step on the squared error $\tfrac12\norm{\mS\vk_t-\vv_t}^2$ of the newest pair.
Both are \emph{heuristic estimators}: they move the memory toward something, but nobody asked
what the best memory would be given everything seen so far, or how sure the memory is of what it
stores.

This chapter takes the question literally and treats memory as \emph{estimation}. Two
consequences follow. First, the best estimate of $\mS$ given all past pairs is a classical object
(the recursive-least-squares solution, or the posterior mean of a Kalman filter), and it can be
computed exactly or to any accuracy we like, instead of being approximated by one gradient step.
Second, an estimator knows how uncertain it is, and that uncertainty can decide how strongly new
evidence is written. We study three 2026 systems that use these ideas in different ways.
\emph{Kalman Linear Attention} (KLA, \citet{shaj2026kalman}; in this chapter KLA always means
the Kalman layer, not the Kaczmarz layer of \cref{ch:delta}) runs an exact Kalman filter inside
every state cell and shows that, written in information form, the filter becomes a parallel
scan. \emph{Gated KalmaNet} (GKA, \citet{peng2026gka}) makes the state the solution of a ridge
regression over the gated, fading past and solves it at every token with Chebyshev iteration.
\emph{Next-Latent Prediction} (NextLat, \citet{teoh2026nextlat}) changes no write rule at all: it
adds a training objective that makes a Transformer's own hidden states behave like the belief
state of a filter, so the memory is shaped offline by the loss.

\Cref{sec:c08-background} builds the estimation toolkit from scratch: Gaussian conditioning, the
Kalman filter and its information form, associative scans, recursive least squares and its link
to the delta rule, ridge regression with forgetting, iterative linear solvers and belief states.
\Cref{sec:c08-kla,sec:c08-gka,sec:c08-nextlat} treat the three systems.
\Cref{sec:c08-comparison} compares them and \cref{sec:c08-estimation} closes with the unifying
``memory as estimation'' view, in which the delta rule, RLS, the Kalman filter and GKA are points
on one line from a single first-order step to an exact second-order solve.

% ===========================================================================================
\section{Background: memory as estimation}\label{sec:c08-background}

\subsection{Recap: linear attention and the delta rule}\label{sec:c08-recap}

We use the notation of \cref{ch:linear,ch:delta}. A head receives per-token keys
$\vk_t\in\R^{d_k}$, values $\vv_t\in\R^{d_v}$ and queries $\vq_t\in\R^{d_k}$, keeps a matrix state
$\mS_t\in\R^{d_v\times d_k}$ and reads $\vo_t=\mS_t\vq_t$. Linear attention writes
\begin{equation}\label{eq:c08-la}
  \mS_t = \mS_{t-1} + \vv_t\vk_t\T ,
\end{equation}
and gated linear attention (GLA, \citet{yang2024gla}) multiplies the old state by a decay first,
$\mS_t = \alpha_t\mS_{t-1}+\vv_t\vk_t\T$. The delta rule \citep{c02-schlag2021linear,yang2025gdn}
writes
\begin{equation}\label{eq:c08-delta}
  \mS_t = \mS_{t-1} - \beta_t(\mS_{t-1}\vk_t-\vv_t)\vk_t\T
        = \mS_{t-1}(\mI-\beta_t\vk_t\vk_t\T)+\beta_t\vv_t\vk_t\T .
\end{equation}
The gradient of $\ell_t(\mS)=\tfrac12\norm{\mS\vk_t-\vv_t}^2$ with respect to $\mS$ is
$(\mS\vk_t-\vv_t)\vk_t\T$, so \eqref{eq:c08-delta} is exactly one gradient-descent step with
learning rate $\beta_t$ on the loss of the newest pair, started from $\mS_{t-1}$. Linear attention
\eqref{eq:c08-la} is likewise one unit step on the ``correlation'' loss $-\inner{\mS\vk_t}{\vv_t}$,
whose gradient is $-\vv_t\vk_t\T$.

The regression problem a key--value memory is really trying to solve involves \emph{all} pairs:
\begin{equation}\label{eq:c08-ls}
  \loss_t(\mS) = \frac12\sum_{i=1}^{t}\norm{\mS\vk_i-\vv_i}^2 + \frac{\lambda}{2}\norm{\mS}_F^2 .
\end{equation}
One gradient step on the newest term $\ell_t$ does not minimize $\loss_t$; it can even undo what
earlier steps stored. The rest of this section builds the tools to minimize \eqref{eq:c08-ls}
exactly, recursively and in parallel, and to reason about uncertainty while doing so.

\subsection{Two Gaussian facts}\label{sec:c08-gauss}

A Gaussian $\vz\sim\mathcal{N}(\vm,\mP)$ in $\R^n$ has density proportional to
$\exp\bigl(-\tfrac12(\vz-\vm)\T\mP^{-1}(\vz-\vm)\bigr)$. The inverse covariance
$\mLambda=\mP^{-1}$ is the \emph{precision}. Two facts carry the whole chapter.

\paragraph{Fact 1 (linear maps).} If $\vz\sim\mathcal{N}(\vm,\mP)$ and $\vw\sim\mathcal{N}(\bm{0},
\mSigma_w)$ is independent of $\vz$, then $\mA\vz+\vw\sim\mathcal{N}(\mA\vm,\mA\mP\mA\T+\mSigma_w)$.
The mean follows from linearity of expectation; the covariance from
$\E[(\mA(\vz-\vm)+\vw)(\mA(\vz-\vm)+\vw)\T]=\mA\mP\mA\T+\mSigma_w$, the cross terms vanishing by
independence. A linear function of jointly Gaussian variables is Gaussian.

\paragraph{Fact 2 (conditioning).} Let the prior be $\vz\sim\mathcal{N}(\vm^-,\mP^-)$ and let the
observation be $\vy=\mC\vz+\ve$ with $\ve\sim\mathcal{N}(\bm{0},\mSigma_e)$. By Bayes' rule the
log-posterior is, up to a constant,
\begin{align*}
  \log p(\vz\mid\vy) &= -\tfrac12(\vz-\vm^-)\T(\mP^-)^{-1}(\vz-\vm^-)
                         -\tfrac12(\vy-\mC\vz)\T\mSigma_e^{-1}(\vy-\mC\vz) \\
  &= -\tfrac12\vz\T\Bigl[(\mP^-)^{-1}+\mC\T\mSigma_e^{-1}\mC\Bigr]\vz
     + \vz\T\Bigl[(\mP^-)^{-1}\vm^-+\mC\T\mSigma_e^{-1}\vy\Bigr] + \text{const}.
\end{align*}
A log-density that is quadratic in $\vz$ belongs to a Gaussian, and comparing with
$-\tfrac12\vz\T\mP^{-1}\vz+\vz\T\mP^{-1}\vm$ (the expansion of the Gaussian exponent) gives the
posterior $\mathcal{N}(\vm,\mP)$ with
\begin{equation}\label{eq:c08-cond-info}
  \mP^{-1} = (\mP^-)^{-1} + \mC\T\mSigma_e^{-1}\mC, \qquad
  \mP^{-1}\vm = (\mP^-)^{-1}\vm^- + \mC\T\mSigma_e^{-1}\vy .
\end{equation}
Notice what happened: \emph{precisions add}, and so do the vectors $\mP^{-1}\vm$. We will exploit
this additivity repeatedly.

\subsection{Linear-Gaussian state-space models and the Kalman filter}\label{sec:c08-kf}

A linear-Gaussian state-space model has a hidden state $\vz_t\in\R^n$ and observations
$\vy_t$:
\begin{equation}\label{eq:c08-lgssm}
  \vz_t = \mA\vz_{t-1}+\vw_t,\quad \vw_t\sim\mathcal{N}(\bm{0},\mSigma_w);\qquad
  \vy_t = \mC_t\vz_t+\ve_t,\quad \ve_t\sim\mathcal{N}(\bm{0},\mSigma_{e,t}),
\end{equation}
with $\vz_0\sim\mathcal{N}(\vm_0,\mP_0)$ and all noises independent. $\mA$ is the transition,
$\mSigma_w$ the \emph{process noise} (how much the state drifts per step) and $\mSigma_{e,t}$ the
\emph{observation noise} (how unreliable each measurement is). The \emph{filtering} problem asks
for $p(\vz_t\mid\vy_{1:t})$. Because $\vz_t$ is Markov and $\vy_t$ depends only on $\vz_t$, the
filtering distribution obeys a two-step recursion: \emph{predict}
$p(\vz_t\mid\vy_{1:t-1})=\int p(\vz_t\mid\vz_{t-1})\,p(\vz_{t-1}\mid\vy_{1:t-1})\,d\vz_{t-1}$, then
\emph{update} $p(\vz_t\mid\vy_{1:t})\propto p(\vy_t\mid\vz_t)\,p(\vz_t\mid\vy_{1:t-1})$. If the
filtering distribution at $t-1$ is Gaussian, Fact 1 makes the prediction Gaussian and Fact 2 makes
the update Gaussian, so by induction everything stays Gaussian and we only need to track a mean
and a covariance. This recursion is the Kalman filter \citep{c08-kalman1960,c08-sarkka2013bayesian}.

\paragraph{Predict.} By Fact 1, with $\vm_t^-$, $\mP_t^-$ the predicted (prior) moments,
\begin{equation}\label{eq:c08-kf-predict}
  \vm_t^- = \mA\vm_{t-1}, \qquad \mP_t^- = \mA\mP_{t-1}\mA\T+\mSigma_w .
\end{equation}

\paragraph{Update.} Fact 2 gives the update in the precision form \eqref{eq:c08-cond-info}. The
textbook \emph{covariance form} avoids inverting state-sized $n\times n$ matrices; it inverts only
an observation-sized one. Define the innovation
$\vy_t-\mC_t\vm_t^-$, its covariance $\bm{E}_t=\mC_t\mP_t^-\mC_t\T+\mSigma_{e,t}$ and the
\emph{Kalman gain} (we write $\mG$ because $\mK$ already denotes stacked keys)
\begin{equation}\label{eq:c08-gain}
  \mG_t = \mP_t^-\mC_t\T\bm{E}_t^{-1}.
\end{equation}
Then
\begin{equation}\label{eq:c08-kf-update}
  \vm_t = \vm_t^- + \mG_t(\vy_t-\mC_t\vm_t^-), \qquad \mP_t = (\mI-\mG_t\mC_t)\mP_t^- .
\end{equation}
\emph{Derivation from \eqref{eq:c08-cond-info}.} The Woodbury identity
$(\mP_-^{-1}+\mC\T\mSigma_e^{-1}\mC)^{-1}=\mP^- -\mP^-\mC\T(\mC\mP^-\mC\T+\mSigma_e)^{-1}\mC\mP^-$
(verify by multiplying out; we write $\mP_-^{-1}$ for $(\mP^-)^{-1}$) gives $\mP=\mP^- - \mG\mC\mP^-$, the covariance update. For the mean
we first show the \emph{gain identity} $\mP\mC\T\mSigma_e^{-1}=\mG$:
\[
  \mP\mC\T\mSigma_e^{-1}
  = \mP^-\mC\T\bigl(\mI-\bm{E}^{-1}\mC\mP^-\mC\T\bigr)\mSigma_e^{-1}
  = \mP^-\mC\T\bm{E}^{-1}\bigl(\bm{E}-\mC\mP^-\mC\T\bigr)\mSigma_e^{-1}
  = \mP^-\mC\T\bm{E}^{-1} .
\]
Multiplying the second equation of \eqref{eq:c08-cond-info} by $\mP$ and using
$\mP\mP_-^{-1}=\mI-\mP\mC\T\mSigma_e^{-1}\mC$ (rearrange the first equation) gives
$\vm=\vm^-+\mP\mC\T\mSigma_e^{-1}(\vy-\mC\vm^-)=\vm^-+\mG(\vy-\mC\vm^-)$.

\paragraph{What the gain means.} Take a scalar state with $a=1$, no process noise, $c=1$ and
observation variance $\sigma^2$. The gain is $G_t=P_t^-/(P_t^-+\sigma^2)$: if the memory is very
uncertain ($P^-\gg\sigma^2$) it copies the observation; if it is very sure ($P^-\ll\sigma^2$) it
barely moves. With $P_0=\sigma^2=1$ the recursion gives $P_t=1/(t+1)$ and $G_t=1/(t+1)$, so
$m_t=\frac{1}{t+1}\sum_{i\le t}y_i$: the filter computes a running mean whose \emph{learning rate
decays as evidence accumulates}. A delta rule with fixed $\beta$ would instead keep forgetting old
observations at a constant rate. This contrast, an adaptive, uncertainty-driven step size versus a
fixed one, is the theme of the chapter.

\subsection{The information form}\label{sec:c08-info}

Write the filter in terms of the precision $\mLambda_t=\mP_t^{-1}$ and the \emph{information
vector} $\vu_t=\mLambda_t\vm_t$. By \eqref{eq:c08-cond-info} the update becomes
\begin{equation}\label{eq:c08-info-update}
  \mLambda_t = \mLambda_t^- + \mC_t\T\mSigma_{e,t}^{-1}\mC_t,\qquad
  \vu_t = \vu_t^- + \mC_t\T\mSigma_{e,t}^{-1}\vy_t .
\end{equation}
Each observation contributes a fixed packet of information, a matrix and a vector that depend only
on that observation, and the packets simply \emph{add}. The price is paid in the predict step,
which in information form reads
\begin{equation}\label{eq:c08-info-predict}
  \mLambda_t^- = \bigl(\mA\mLambda_{t-1}^{-1}\mA\T+\mSigma_w\bigr)^{-1},\qquad
  \vu_t^- = \mLambda_t^-\mA\mLambda_{t-1}^{-1}\vu_{t-1} .
\end{equation}
The covariance and information forms are dual: one has an easy predict and a hard update, the
other the reverse. Two special cases matter for this chapter.
\begin{itemize}
  \item \textbf{No process noise} ($\mSigma_w=\bm{0}$, $\mA$ invertible). Then
  $\mLambda_t^-=\mA^{-\top}\mLambda_{t-1}\mA^{-1}$ and $\vu_t^-=\mA^{-\top}\vu_{t-1}$, both linear.
  With a static state ($\mA=\mI$) the precision is just a running sum of information packets,
  exactly like the linear-attention state \eqref{eq:c08-la}.
  \item \textbf{Scalar state with process noise.} With $\mA=a$, $\mSigma_w=p$ and precision
  $\Lambda$, \eqref{eq:c08-info-predict} becomes
  $\Lambda_t^-=\bigl(a^2/\Lambda_{t-1}+p\bigr)^{-1}=\Lambda_{t-1}/(a^2+p\Lambda_{t-1})$, a
  \emph{linear-fractional} (M\"obius) function of $\Lambda_{t-1}$. It is not linear, but, as we
  show next, linear-fractional maps compose as easily as linear ones.
\end{itemize}
Why does additivity matter? A recurrence whose step is ``apply a fixed map, then add a fixed
term'' can be evaluated for all $t$ at once with a parallel scan, which is how linear attention is
trained (\cref{ch:linear}). The covariance-form gain \eqref{eq:c08-gain} hides this structure
behind a matrix inverse that depends on the running state.

\subsection{Associative scans}\label{sec:c08-scan}

Given elements $e_1,\dots,e_T$ and an associative operation $\oplus$, the \emph{inclusive scan}
returns all prefixes $s_t=e_1\oplus e_2\oplus\dots\oplus e_t$. Associativity lets us bracket the
products as a balanced tree, so all prefixes can be computed in $\bigO(\log T)$ parallel steps.
The Hillis--Steele scheme \citep{c08-hillis1986data} does $\bigO(T\log T)$ work: at round
$r=0,1,\dots$ every position $t$ with $t>2^r$ combines its value with the value $2^r$ positions to
its left. The Blelloch scheme \citep{c03-blelloch1990prefix} does $\bigO(T)$ work with an up-sweep and
a down-sweep. Both need only the combine rule. Two combine rules appear in this chapter.

\paragraph{Affine maps.} The recurrence $z_t=a_tz_{t-1}+b_t$ applies the map
$f_t(z)=a_tz+b_t$. Composition of affine maps is affine:
\[
  f_2\circ f_1(z) = a_2(a_1z+b_1)+b_2 = (a_2a_1)z+(a_2b_1+b_2),
\]
so the pairs combine as $(a_1,b_1)\oplus(a_2,b_2)=(a_2a_1,\;a_2b_1+b_2)$, and $z_t$ is the second
component of the prefix applied to $z_0=0$. Gated linear attention is this recurrence with
matrix-valued $b_t=\vv_t\vk_t\T$.

\paragraph{M\"obius maps.} A map $g(\lambda)=(A\lambda+B)/(C\lambda+D)$ is identified with the
matrix $\left[\begin{smallmatrix}A&B\\C&D\end{smallmatrix}\right]$. Substituting one into another,
\[
  g_2(g_1(\lambda)) = \frac{A_2\frac{A_1\lambda+B_1}{C_1\lambda+D_1}+B_2}
                           {C_2\frac{A_1\lambda+B_1}{C_1\lambda+D_1}+D_2}
  = \frac{(A_2A_1+B_2C_1)\lambda+(A_2B_1+B_2D_1)}{(C_2A_1+D_2C_1)\lambda+(C_2B_1+D_2D_1)} ,
\]
whose coefficients are exactly the entries of the matrix product
$\left[\begin{smallmatrix}A_2&B_2\\C_2&D_2\end{smallmatrix}\right]
\left[\begin{smallmatrix}A_1&B_1\\C_1&D_1\end{smallmatrix}\right]$. Matrix multiplication is
associative, so M\"obius recurrences are scans over $2\times2$ matrices. Multiplying a matrix by
any nonzero scalar does not change the map (numerator and denominator scale together), a freedom
we will use for numerical stability. \Citet{c08-sarkka2021temporal} used scans of this kind to
parallelize general Bayesian filters and smoothers over time; KLA brings the idea into a
language-model layer.

\subsection{Recursive least squares}\label{sec:c08-rls}

Minimize the ridge objective \eqref{eq:c08-ls}. Setting its gradient
$\sum_i(\mS\vk_i-\vv_i)\vk_i\T+\lambda\mS$ to zero gives
\begin{equation}\label{eq:c08-ridge}
  \mS_t = \mU_t(\mH_t+\lambda\mI)^{-1},\qquad
  \mH_t=\sum_{i\le t}\vk_i\vk_i\T,\quad \mU_t=\sum_{i\le t}\vv_i\vk_i\T .
\end{equation}
$\mH_t$ is the (uncentred) key covariance and $\mU_t$ is precisely the linear-attention state
\eqref{eq:c08-la}. Solving \eqref{eq:c08-ridge} from scratch costs $\bigO(d_k^3)$ per token.
\emph{Recursive least squares} (RLS) instead propagates $\mP_t=(\mH_t+\lambda\mI)^{-1}$. Since
$\mP_t^{-1}=\mP_{t-1}^{-1}+\vk_t\vk_t\T$ is a rank-one change, the Sherman--Morrison formula gives
\begin{equation}\label{eq:c08-rls-P}
  \mP_t = \mP_{t-1}-\frac{\mP_{t-1}\vk_t\vk_t\T\mP_{t-1}}{1+\vk_t\T\mP_{t-1}\vk_t},
\end{equation}
which costs $\bigO(d_k^2)$. For the state, write $\mU_{t-1}=\mS_{t-1}\mP_{t-1}^{-1}=
\mS_{t-1}(\mP_t^{-1}-\vk_t\vk_t\T)$, so
\begin{equation}\label{eq:c08-rls-S}
  \mS_t = (\mU_{t-1}+\vv_t\vk_t\T)\mP_t
        = \mS_{t-1} + (\vv_t-\mS_{t-1}\vk_t)\,\vk_t\T\mP_t .
\end{equation}

\paragraph{RLS is a Kalman filter.} Treat each row $\vs_j\T$ of $\mS$ as a state vector that never
changes ($\mA=\mI$, $\mSigma_w=\bm{0}$), give it the prior $\mathcal{N}(\bm{0},\lambda^{-1}\mI)$, and
observe $v_t[j]=\vk_t\T\vs_j+e$ with unit noise. The information form
\eqref{eq:c08-info-update} gives $\mLambda_t=\lambda\mI+\sum_i\vk_i\vk_i\T=\mH_t+\lambda\mI$ and,
stacking the rows, the information matrix $\mU_t$; the posterior mean is
$\mU_t\mLambda_t^{-1}$, which is \eqref{eq:c08-ridge}. The Kalman gain for this model is
$\mP_t^-\vk_t/(\vk_t\T\mP_t^-\vk_t+1)=\mP_t\vk_t$ (the gain identity), and \eqref{eq:c08-kf-update}
reproduces \eqref{eq:c08-rls-S}. So the ridge memory, the RLS memory and the posterior mean of a
static Kalman filter are the same object.

\paragraph{RLS is a Newton step.} The Hessian of $\loss_t$ in \eqref{eq:c08-ls} (acting on rows of
$\mS$) is $\mH_t+\lambda\mI=\mP_t^{-1}$. Since $\mS_{t-1}$ minimizes $\loss_{t-1}$, the gradient of
$\loss_t=\loss_{t-1}+\ell_t$ at $\mS_{t-1}$ is just $\nabla\ell_t(\mS_{t-1})=(\mS_{t-1}\vk_t-\vv_t)
\vk_t\T$. A Newton step from $\mS_{t-1}$ is therefore
$\mS_{t-1}-(\mS_{t-1}\vk_t-\vv_t)\vk_t\T\mP_t$, which is \eqref{eq:c08-rls-S}; because $\loss_t$ is
quadratic, one Newton step lands exactly on its minimizer.

\begin{keyidea}[title={Key idea: the delta rule is RLS with an identity covariance}]
Compare the RLS update \eqref{eq:c08-rls-S} with the delta rule \eqref{eq:c08-delta}:
\[
  \text{RLS: }\ \mS_t=\mS_{t-1}+(\vv_t-\mS_{t-1}\vk_t)\vk_t\T\,\mP_t, \qquad
  \text{delta: }\ \mS_t=\mS_{t-1}+(\vv_t-\mS_{t-1}\vk_t)\vk_t\T\,\beta_t\mI .
\]
The delta rule is RLS in which the inverse Hessian $\mP_t$ (the posterior covariance) is replaced
by the isotropic, constant $\beta_t\mI$. RLS gives each key direction its own learning rate:
directions seen often have small posterior variance and are updated gently; rare directions are
updated strongly. The delta rule forgets this bookkeeping and pays for it with interference.
\end{keyidea}

\subsection{Ridge regression with exponential forgetting}\label{sec:c08-ewridge}

A memory for language must forget. Weight pair $i$ at time $t$ by $w_{t,i}=\prod_{j=i+1}^{t}
\alpha_j$ with gates $\alpha_j\in(0,1]$ (a constant $\gamma$ gives $w_{t,i}=\gamma^{t-i}$):
\begin{equation}\label{eq:c08-ewls}
  \loss_t(\mS)=\frac12\sum_{i\le t}w_{t,i}\norm{\mS\vk_i-\vv_i}^2+\frac{\lambda}{2}\norm{\mS}_F^2 ,
  \qquad
  \mS_t=\mU_t(\mH_t+\lambda\mI)^{-1},
\end{equation}
where now $\mH_t=\alpha_t\mH_{t-1}+\vk_t\vk_t\T$ and $\mU_t=\alpha_t\mU_{t-1}+\vv_t\vk_t\T$. Both
are gated-linear-attention states: $\mU_t$ with values $\vv$, $\mH_t$ with values $\vk$. A ridge
memory is therefore ``two GLA memories plus a linear solve at read time''.

There are two ways to place the regularizer. Classical RLS with a forgetting factor $\gamma$ decays
the prior together with the data, which keeps the rank-one recursion
\eqref{eq:c08-rls-P} (use $\mP_{t-1}/\gamma$ in place of $\mP_{t-1}$) but makes the effective
regularizer $\lambda\gamma^t$ vanish. When keys stop exciting some direction, $\mH_t$ loses rank
in that direction and $\mP_t$ grows without bound, a failure known as covariance wind-up. Keeping
$\lambda$ fixed, as in \eqref{eq:c08-ewls}, prevents this, but then
$\mP_t^{-1}=\alpha_t\mP_{t-1}^{-1}+\vk_t\vk_t\T+(1-\alpha_t)\lambda\mI$ is no longer a rank-one
update and Sherman--Morrison no longer applies. This is one reason Gated KalmaNet solves the linear
system iteratively instead of propagating $\mP_t$.

\subsection{Iterative solvers: gradient descent, Chebyshev, conjugate gradient}\label{sec:c08-solvers}

We need to solve $\mA\vx=\vb$ with $\mA$ symmetric positive definite and eigenvalues in a known
interval $[\mu_{\min},\mu_{\max}]$; $\kappa=\mu_{\max}/\mu_{\min}$ is the condition number.

\paragraph{Gradient descent (Richardson iteration).} $\vx_{n+1}=\vx_n-\tau(\mA\vx_n-\vb)$ is
gradient descent on $\tfrac12\vx\T\mA\vx-\vb\T\vx$. The error $\ve_n=\vx_n-\vx^\star$ obeys
$\ve_{n+1}=(\mI-\tau\mA)\ve_n$, so each eigencomponent is multiplied by $1-\tau\mu$. The step
$\tau=2/(\mu_{\min}+\mu_{\max})$ balances the two ends of the spectrum and gives the contraction
factor $\rho=(\mu_{\max}-\mu_{\min})/(\mu_{\max}+\mu_{\min})=(\kappa-1)/(\kappa+1)$; reaching
relative error $\epsilon$ takes about $\tfrac{\kappa}{2}\ln(1/\epsilon)$ iterations.

\paragraph{The polynomial view.} Any method whose iterate $\vx_n$ lies in
$\mathrm{span}\{\vb,\mA\vb,\dots,\mA^{n-1}\vb\}$ (started at $\vx_0=\bm{0}$) has error
$\ve_n=p_n(\mA)\ve_0$ for a polynomial $p_n$ of degree $n$ with $p_n(0)=1$. Richardson uses
$p_n(\mu)=(1-\tau\mu)^n$. The best worst-case choice minimizes $\max_{\mu\in[\mu_{\min},\mu_{\max}]}
\abs{p_n(\mu)}$, and the answer is a rescaled Chebyshev polynomial. The Chebyshev polynomials are
defined by $T_0(\xi)=1$, $T_1(\xi)=\xi$ and
\begin{equation}\label{eq:c08-cheb-rec}
  T_{n+1}(\xi)=2\xi T_n(\xi)-T_{n-1}(\xi);
\end{equation}
they satisfy $\abs{T_n(\xi)}\le1$ on $[-1,1]$ and grow as fast as possible outside it, with
$T_n(\xi)=\tfrac12\bigl[(\xi+\sqrt{\xi^2-1})^n+(\xi-\sqrt{\xi^2-1})^n\bigr]$ for $\xi\ge1$. Map the
spectrum to $[-1,1]$ with $\xi(\mu)=(1-\tau\mu)/\rho$ and set
\begin{equation}\label{eq:c08-cheb-poly}
  p_n(\mu)=\frac{T_n(\xi(\mu))}{T_n(1/\rho)},\qquad\text{so}\qquad
  \max_{\mu}\abs{p_n(\mu)}=\frac{1}{T_n(1/\rho)}\le 2\sigma^n,\quad
  \sigma=\frac{\sqrt\kappa-1}{\sqrt\kappa+1}.
\end{equation}
The bound follows from $T_n(1/\rho)\ge\tfrac12(1/\rho+\sqrt{1/\rho^2-1})^n$ and, with
$1/\rho=(\kappa+1)/(\kappa-1)$, from
$1/\rho+\sqrt{1/\rho^2-1}=(\kappa+1+2\sqrt\kappa)/(\kappa-1)=(\sqrt\kappa+1)/(\sqrt\kappa-1)=1/\sigma$.
Chebyshev needs about $\tfrac{\sqrt\kappa}{2}\ln(2/\epsilon)$ iterations: the square root of what
gradient descent needs.

\paragraph{The three-term recurrence.} We never form $p_n$. Let $c_n=T_n(1/\rho)$ and substitute
$\xi=(1-\tau\mu)/\rho$ into \eqref{eq:c08-cheb-rec}: $c_{n+1}p_{n+1}=\tfrac{2}{\rho}(1-\tau\mu)c_np_n
-c_{n-1}p_{n-1}$. Dividing by $c_{n+1}$ and defining $\omega_{n+1}=2c_n/(\rho c_{n+1})$, the
recurrence $c_{n+1}=\tfrac2\rho c_n-c_{n-1}$ gives $c_{n-1}/c_{n+1}=\omega_{n+1}-1$ and
\[
  \frac{1}{\omega_{n+1}}=\frac{\rho c_{n+1}}{2c_n}=1-\frac{\rho c_{n-1}}{2c_n}
  =1-\frac{\rho^2}{4}\,\omega_n .
\]
Hence $p_{n+1}=\omega_{n+1}(1-\tau\mu)p_n+(1-\omega_{n+1})p_{n-1}$, which, applied to the error,
is the iteration \citep{c08-golub1961chebyshev}
\begin{equation}\label{eq:c08-cheb-iter}
  \vx_{n+1}=\omega_{n+1}\bigl(\vx_n-\tau(\mA\vx_n-\vb)\bigr)+(1-\omega_{n+1})\vx_{n-1},\qquad
  \omega_{n+1}=\Bigl(1-\tfrac{\rho^2}{4}\omega_n\Bigr)^{-1},
\end{equation}
started from $\vx_0=\bm{0}$, $\vx_1=\tau\vb$ (one Richardson step, $p_1(\mu)=1-\tau\mu$) and
$\omega_1=2c_0/(\rho c_1)=2$. Each iteration costs one matrix--vector product, like gradient
descent; the weights $\omega_n$ depend only on the interval, never on the data.

\paragraph{Conjugate gradient.} CG \citep{c08-hestenes1952cg} picks the best polynomial for the
actual spectrum (it minimizes the $\mA$-norm error over the same Krylov space), so it satisfies the
same $2\sigma^n$ bound (in the $\mA$-norm) and is often much faster. But its step sizes are ratios of inner products of
computed residuals. That has three costs on accelerators: two global reductions per iteration;
coefficients that drift when residuals lose orthogonality in low precision; and an output that is a
\emph{nonlinear} function of $\vb$. Chebyshev's output $\vx_n=q_n(\mA)\vb$ is a fixed polynomial
in $\mA$ applied to $\vb$: linear in $\vb$, with a symmetric Jacobian $q_n(\mA)$, so its backward
pass is the same solver applied to the incoming gradient. GKA exploits exactly this.

\begin{example}[Solver race]\label{ex:c08-race}
In our NumPy check we built $\mH$ as a decayed sum of 500 unit-norm keys in $d=64$ dimensions and
solved $(\mH+\lambda\mI)\vx=\vq$ with $\lambda=0.02\norm{\mH}_F$ (the GKA recipe of
\cref{sec:c08-gka}); the true condition number was 33.7. After 30 iterations the relative error was
$1.1\times10^{-1}$ for optimally tuned gradient descent, $2.4\times10^{-4}$ for Chebyshev (the
bound $2\sigma^{30}$ with $\kappa\le51$ is $4.2\times10^{-4}$) and at machine precision for CG in
float64. When we rounded every operation to float16, both Chebyshev and CG stalled at the float16
floor of about $10^{-3}$, so this small experiment does not reproduce the low-precision instability
that motivates GKA; the paper reports that comparison at scale.
\end{example}

\subsection{Belief states and world models}\label{sec:c08-belief}

Consider a hidden Markov model: a hidden state $s_t$ evolves by $P(s_{t+1}\mid s_t)$ and emits a
token $x_t\sim P(x_t\mid s_t)$. The \emph{belief} $\vb_t(s)=P(s_{t+1}=s\mid x_{1:t})$ is the
predictive distribution of the next hidden state given the history. It obeys a recursion: condition
on the new token, then predict,
\begin{equation}\label{eq:c08-bayes-filter}
  \vb_t = f(\vb_{t-1},x_t),\qquad
  f(\vb,x)(s') = \sum_{s}P(s'\mid s)\,\frac{\vb(s)P(x\mid s)}{\sum_{\tilde s}\vb(\tilde s)P(x\mid\tilde s)} .
\end{equation}
The Kalman filter is this \emph{Bayes filter} for linear-Gaussian models, where the belief is the
pair $(\vm_t,\mP_t)$. The crucial property is \emph{sufficiency}: for every horizon $k$,
\[
  P(x_{t+1:t+k}\mid x_{1:t}) = \prod_{j=1}^{k}\Bigl(\sum_s \vb_{t+j-1}(s)P(x_{t+j}\mid s)\Bigr),
  \qquad \vb_{t+j}=f(\vb_{t+j-1},x_{t+j}),
\]
so the future depends on the past only through $\vb_t$. A belief state is the ideal memory: as
small as the task allows, yet sufficient, and updatable by a fixed rule. In POMDP terms
\citep{c08-kaelbling1998pomdp} it is the state on which optimal decisions depend. A pair
(transition rule $f$, read-out from $\vb$ to the next-token distribution) is a \emph{world model}.
In our NumPy check, rolling the belief of a random 4-state HMM forward with its own update predicted
the joint probability of every 1-, 2- and 3-token continuation to within $6\times10^{-17}$ of a
brute-force computation over the full history.

Different architectures relate to belief states differently. An RNN or SSM must compress the
history into a fixed state, but nothing forces that state to be sufficient. A Transformer keeps the
whole history in its KV cache and can look anything up later, so it has no incentive to compress at
all. KLA computes an exact belief for a simple factorized model; GKA computes the posterior mean of
a static regression model; NextLat trains a Transformer so that its hidden state behaves like a
belief state of the data.

% ===========================================================================================
\section{Kalman Linear Attention}\label{sec:c08-kla}

\subsection{Motivation}

In Mamba and gated linear attention the state is a \emph{point estimate}. The decay gate
$\alpha_t$ is computed from the current token alone; it never asks how much the memory already
knows. In a Kalman filter, by contrast, the amount of new evidence written into the state, the
gain, is a function of the state's own uncertainty: an uncertain memory listens, a confident one
does not. \Citet{shaj2026kalman} propose to make this mechanism the sequence mixer of a language
model. Kalman Linear Attention keeps a Gaussian \emph{belief} for every state cell, derives the
write gate from that belief, and keeps parallel training by working in the information form of
\cref{sec:c08-info}, where the filter becomes an associative scan. The repository README
summarizes the pitch: ``Unlike other linear attention layers which model the current state as a
single point in the state space, KLA models the current state as a belief over the state space,
and updates it in closed form as the sequence arrives''
(\repofile{repos/vaisakh-shaj__kalman-linear-attention/readme.md}).

\begin{taxonomy}
Following \citet{zhoubian2026memory}, KLA is \textbf{implicit} memory: the belief lives inside the
layer and is read and written only by the forward computation. Its \textbf{update dynamics} are
\textbf{online} (the belief is revised at every token, at inference as well as in training) and its
\textbf{persistence} is \textbf{long-term} within a sequence (survey Table~I lists Kalman Linear
Attention and Gated KalmaNet together as ``Online, Long-Term, Probabilistic Recurrent State
Estimation''). Its \textbf{update rule} (survey Table~II) is a \emph{state-transition update} given
by a filtering equation; the survey names Kalman Linear Attention as a representative example and
describes Kalman-style variants as introducing ``uncertainty-aware filtering''.
\end{taxonomy}

\subsection{Formulation: a Kalman filter in every cell}\label{sec:c08-kla-form}

KLA's layout follows Mamba (\cref{ch:linear}). The layer has $d_v$ channels (the code's
\code{d_inner}, $M$) and $d_k$ states per channel (the code's \code{d_state}, $S$, default 16); the
state is a $d_v\times d_k$ array of cells $(m,s)$. Each token supplies a value $\vv_t\in\R^{d_v}$, a
positive \emph{value precision} $\bm{\pi}_t\in\R^{d_v}$ (how reliable the token's value is), and a
key $\vk_t\in\R^{d_k}$ and query $\vq_t\in\R^{d_k}$ shared by all channels, like Mamba's $B_t$ and
$C_t$. Every cell also has a learned decay $a^{ms}\in(0,1)$ and process noise $p^{ms}>0$, both
time-invariant.

The docstring of the ops module,
\repofile{repos/vaisakh-shaj__kalman-linear-attention/src/kla/ops/kla_ops.py}, describes the core
operation as ``a parallel Kalman filter in information form, factorized per (channel m, state s)
pair''.
Reading the update off the code, each cell runs the scalar filter of the model
\begin{equation}\label{eq:c08-kla-model}
  z_t^{ms} = a^{ms}z_{t-1}^{ms}+w_t,\quad w_t\sim\mathcal{N}(0,p^{ms});\qquad
  v_t[m] = k_t[s]\,z_t^{ms}+e_t,\quad e_t\sim\mathcal{N}\bigl(0,1/\pi_t[m]\bigr),
\end{equation}
in which every cell of channel $m$ treats the token's value $v_t[m]$ as its own noisy measurement,
seen through the scalar observation coefficient $k_t[s]$. The cells are filtered independently (a
mean-field factorization), which is what keeps the cost at $\bigO(d_vd_k)$ per token. This is our
reading of the implementation; the paper's own derivation may organize it differently.

\paragraph{One cell, step by step.} Drop the superscripts and let $\Lambda$ be the cell's
precision and $u=\Lambda m$ its information. The predict step \eqref{eq:c08-info-predict} for a
scalar state gives $\Lambda_t^-=\Lambda_{t-1}/(a^2+p\Lambda_{t-1})$ and
$u_t^-=\Lambda_t^-\,a\,(u_{t-1}/\Lambda_{t-1})=a\,u_{t-1}/(a^2+p\Lambda_{t-1})$. The update
\eqref{eq:c08-info-update} with $c=k_t[s]$, noise variance $1/\pi_t[m]$ and measurement $v_t[m]$
adds the information packet
\begin{equation}\label{eq:c08-kla-stats}
  \phi_t = \pi_t[m]\,k_t[s]^2,\qquad r_t=\pi_t[m]\,v_t[m]\,k_t[s].
\end{equation}
Altogether,
\begin{align}
  \Lambda_t &= \frac{\Lambda_{t-1}}{a^2+p\Lambda_{t-1}}+\phi_t, \label{eq:c08-kla-prec}\\
  u_t &= \alpha_t\,u_{t-1}+r_t,\qquad \alpha_t=\frac{a}{a^2+p\Lambda_{t-1}}, \label{eq:c08-kla-info}\\
  \text{mean } &= u_t/\Lambda_t,\qquad \text{variance } = 1/\Lambda_t . \nonumber
\end{align}
Stacking the cells into matrices $\mLambda_t,\mU_t\in\R^{d_v\times d_k}$ and the gates into
$\bm{\alpha}_t$, and writing $\oslash$ for element-wise division, the layer is
\begin{equation}\label{eq:c08-kla-matrix}
  \mU_t=\bm{\alpha}_t\had\mU_{t-1}+(\bm{\pi}_t\had\vv_t)\vk_t\T,\qquad
  \vo_t=(\mU_t\oslash\mLambda_t)\,\vq_t,\qquad
  \vo^{\mathrm{var}}_t=(\bm{1}\oslash\mLambda_t)(\vq_t\had\vq_t).
\end{equation}
The first equation is a gated-linear-attention write with value $\bm{\pi}_t\had\vv_t$; the read-out
is $\vo_t=\mS_t\vq_t$ with the posterior mean $\mS_t=\mU_t\oslash\mLambda_t$ playing the role of the
memory matrix. Three things distinguish KLA from GLA:
\begin{enumerate}
  \item the gate $\bm{\alpha}_t$ is computed from the precision, i.e.\ from how much information the
  cell has accumulated, rather than from the current token;
  \item the read-out divides each cell by its precision, a per-cell normalizer that plays the role
  linear attention's denominator $\sum_i\vk_i\T\vq$ plays for the whole state;
  \item each token's write is weighted by its learned precision $\bm{\pi}_t$, so the model can
  declare some tokens unreliable.
\end{enumerate}
The second output $\vo^{\mathrm{var}}_t$ is the variance of $\vo_t$ under the factorized posterior.
This is the ``uncertainty carried in the state'': it costs nothing extra because $1/\Lambda_t$ is
already needed for the mean.

\paragraph{How the gate behaves.} Write the posterior mean of one cell as
\[
  \text{mean}_t=\frac{\Lambda_t^-\,(a\,\text{mean}_{t-1})+\phi_t\,(v_t[m]/k_t[s])}{\Lambda_t^-+\phi_t}
  \qquad(k_t[s]\neq0),
\]
a precision-weighted average of the predicted mean and the ``pseudo-observation''
$v_t[m]/k_t[s]$. Because $\Lambda_t^-=\Lambda_{t-1}/(a^2+p\Lambda_{t-1})<1/p$, process noise caps
the memory's confidence: however much evidence it has seen, the weight on new evidence stays above
$\phi_t/(1/p+\phi_t)$, so the cell never stops listening. Process noise is therefore KLA's
forgetting mechanism. Without it ($p=0$), \eqref{eq:c08-kla-prec} becomes
$\Lambda_t=\Lambda_{t-1}/a^2+\phi_t$, which grows like $a^{-2t}$: in our check, with $a=0.95$ and
$\phi=0.5$ the precision passed $4\times10^{9}$ after 200 steps, and the cell effectively stopped
writing. With constant information $\phi$ and $p>0$, the precision converges to the positive root
of $p\Lambda^2+(a^2-1-p\phi)\Lambda-a^2\phi=0$ (set $\Lambda_t=\Lambda_{t-1}$ in
\eqref{eq:c08-kla-prec} and clear the denominator). For $a=0.95$, $p=0.05$, $\phi=0.5$ this gives
$\Lambda^\star\approx4.47$ and a steady gate $\alpha^\star=a/(a^2+p\Lambda^\star)\approx0.844$.
In general the fixed-point equation $\Lambda^\star(1-1/(a^2+p\Lambda^\star))=\phi>0$ forces
$a^2+p\Lambda^\star>1$, hence $\alpha^\star<a$: in steady state the information decays faster than
the dynamics alone would suggest, because process noise keeps eroding old evidence. In
this regime KLA behaves like GLA with constant decay $0.844$ and a read-out scaled by
$1/\Lambda^\star$; away from it, the gate tracks the information actually received.

\subsection{Parallel form: a M\"obius scan and an affine scan}\label{sec:c08-kla-scan}

The precision recurrence \eqref{eq:c08-kla-prec} is linear-fractional. Clearing the denominator,
$\Lambda_t=\bigl((1+p\phi_t)\Lambda_{t-1}+a^2\phi_t\bigr)/\bigl(p\Lambda_{t-1}+a^2\bigr)$, and
dividing numerator and denominator by $a^2$ gives the leaf matrix
\begin{equation}\label{eq:c08-kla-leaf}
  \Lambda_t = \mathcal{M}_t(\Lambda_{t-1}),\qquad
  \mathcal{M}_t=\begin{bmatrix}(1+p\phi_t)/a^2 & \phi_t\\ p/a^2 & 1\end{bmatrix}.
\end{equation}
By \cref{sec:c08-scan}, $\Lambda_t$ is the M\"obius map of the product
$\mathcal{P}_t=\mathcal{M}_t\mathcal{M}_{t-1}\cdots\mathcal{M}_1$ applied to $\Lambda_0$, and all
products come out of one associative scan over $2\times2$ matrices. Once every $\Lambda_t$ is known,
the gates $\alpha_t=a/(a^2+p\Lambda_{t-1})$ are known too, and the information recurrence
\eqref{eq:c08-kla-info} is an affine scan (\cref{sec:c08-scan}). So the exact filter costs two scans.

\begin{proposition}[Parallel equals sequential]\label{prop:c08-kla-par}
For any inputs, the two-scan procedure returns the same $\Lambda_t$, $u_t$ and outputs as running
\eqref{eq:c08-kla-prec}--\eqref{eq:c08-kla-info} token by token, even when every scan node is
rescaled by an arbitrary positive factor.
\end{proposition}
\begin{proof}
By induction, $\Lambda_t=\mathcal{M}_t\circ\dots\circ\mathcal{M}_1(\Lambda_0)$, and composition of
M\"obius maps is the matrix product (\cref{sec:c08-scan}), which is associative, so any bracketing
the scan uses yields the same matrix $\mathcal{P}_t$. A positive rescaling of any intermediate
matrix multiplies the final numerator and denominator by the same factor and cancels. The same
argument with affine composition gives $u_t$.
\end{proof}

\paragraph{Why trace normalization is safe.} Products of many leaves overflow or underflow in
float32, so the implementations divide each combined matrix by its trace $A+D$. Every leaf is
entrywise nonnegative with positive diagonal (as $p,\phi\ge0$), products of such matrices stay so,
and therefore the trace never vanishes. The leaf determinant is
$\frac{1+p\phi}{a^2}-\frac{p\phi}{a^2}=a^{-2}>0$ and determinants multiply, so every node stays
invertible. In our NumPy check the parallel form (Hillis--Steele scans, trace normalization)
matched sequential filtering to $10^{-15}$, and the per-cell information-form filter matched the
textbook covariance-form Kalman filter \eqref{eq:c08-kf-predict}--\eqref{eq:c08-kf-update} to
$10^{-15}$ as well.

\begin{figure}[t]
\centering
\begin{tikzpicture}[
  box/.style={draw,rounded corners=2pt,align=center,font=\small,minimum height=8mm,inner sep=3pt},
  arr/.style={-{Stealth[length=2mm]},thick}]
\node[box] (in) {tokens\\$\vv_t,\bm{\pi}_t,\vk_t$};
\node[box,right=6mm of in] (stats) {packets\\$\phi_t,\ r_t$};
\node[box,right=6mm of stats,fill=blue!6] (mob) {M\"obius scan\\$2\times2$ products};
\node[box,right=6mm of mob] (lam) {$\Lambda_t$, and\\$\alpha_t$ from $\Lambda_{t-1}$};
\node[box,below=10mm of mob,fill=blue!6] (aff) {affine scan\\$u_t=\alpha_tu_{t-1}+r_t$};
\node[box,below=10mm of lam,xshift=8mm] (out) {$\vo_t=(\mU_t\oslash\mLambda_t)\vq_t$\\and $\vo^{\mathrm{var}}_t$};
\draw[arr] (in) -- (stats);
\draw[arr] (stats) -- node[above,font=\scriptsize] {$\phi_t$} (mob);
\draw[arr] (mob) -- (lam);
\draw[arr] (stats.south east) -- node[below left,font=\scriptsize] {$r_t$} (aff.north west);
\draw[arr] (lam.south west) -- node[above left,font=\scriptsize] {$\alpha_t$} (aff.north east);
\draw[arr] (aff) -- (out);
\draw[arr] (lam) -- node[right,font=\scriptsize] {$\Lambda_t$} (out);
\end{tikzpicture}
\caption{KLA computes an exact (factorized) Kalman filter with two parallel scans. The precision
recurrence is a scan over $2\times2$ M\"obius matrices; given the precisions, the gates are known
and the information recurrence is an ordinary affine (GLA-like) scan.}
\label{fig:c08-kla}
\end{figure}

\paragraph{Discretization.} The decay and process noise come from a continuous-time
Ornstein--Uhlenbeck process $dz=-\abs{a_c}z\,dt+\sqrt{p_c}\,dW$ observed every $\Delta$ time units.
Solving it over one interval gives $a=e^{-\abs{a_c}\Delta}$ and
$p=\frac{p_c}{2\abs{a_c}}(1-e^{-2\abs{a_c}\Delta})$ (the stationary variance $p_c/(2\abs{a_c})$
times the fraction not yet reached). In KLA $\Delta$ is a learned \emph{static} parameter per cell,
unlike Mamba's input-dependent $\Delta_t$; the input-dependence of the write enters only through
$\bm{\pi}_t$, $\vk_t$ and $\vv_t$, and through the gate's dependence on the accumulated precision.

\subsection{Algorithms and complexity}

\begin{algorithm}[t]
\caption{KLA prefill / training (parallel)}\label{alg:c08-kla-prefill}
\begin{algorithmic}[1]
\Require $\vv_t,\bm{\pi}_t\in\R^{d_v}$, $\vk_t,\vq_t\in\R^{d_k}$ for $t=1..T$; decay $a$, noise $p$
  $\in\R^{d_v\times d_k}$; prior $(\mLambda_0,\mU_0)$
\Ensure outputs $\vo_t$, variances $\vo^{\mathrm{var}}_t$, final state $(\mLambda_T,\mU_T)$
\State $\phi_t\gets\bm{\pi}_t(\vk_t\had\vk_t)\T$, \ $r_t\gets(\bm{\pi}_t\had\vv_t)\vk_t\T$ \Comment{all $t$, element-wise per cell}
\State leaves $\mathcal{M}_t\gets\bigl[\begin{smallmatrix}(1+p\phi_t)/a^2&\phi_t\\p/a^2&1\end{smallmatrix}\bigr]$
\State $\mathcal{P}_{1..T}\gets$ \Call{Scan}{$\mathcal{M}_{1..T}$, $2\times2$ product, divide by trace}
\State $\mLambda_t\gets(\mathcal{P}^{11}_t\mLambda_0+\mathcal{P}^{12}_t)\oslash(\mathcal{P}^{21}_t\mLambda_0+\mathcal{P}^{22}_t)$
\State $\bm{\alpha}_t\gets a\oslash(a^2+p\had\mLambda_{t-1})$ \Comment{shifted precisions}
\State $\mU_{1..T}\gets$ \Call{Scan}{$(\bm{\alpha}_t,r_t)_{1..T}$, affine composition} applied to $\mU_0$
\State $\vo_t\gets(\mU_t\oslash\mLambda_t)\vq_t$, \ $\vo^{\mathrm{var}}_t\gets(\bm{1}\oslash\mLambda_t)(\vq_t\had\vq_t)$
\end{algorithmic}
\end{algorithm}

\begin{algorithm}[t]
\caption{KLA decode (one token)}\label{alg:c08-kla-decode}
\begin{algorithmic}[1]
\Require state $(\mLambda,\mU)$; new $\vv,\bm{\pi},\vk,\vq$
\State $\bm{D}\gets a^2+p\had\mLambda$
\State $\mLambda\gets\mLambda\oslash\bm{D}+\bm{\pi}(\vk\had\vk)\T$
\State $\mU\gets(a\oslash\bm{D})\had\mU+(\bm{\pi}\had\vv)\vk\T$
\State \Return $\vo=(\mU\oslash\mLambda)\vq$ and the new state
\end{algorithmic}
\end{algorithm}

\Cref{alg:c08-kla-prefill,alg:c08-kla-decode} summarize the layer. The state is two
$d_v\times d_k$ arrays, $2d_vd_k$ numbers, twice that of GLA or Mamba with the same dimensions.
Decoding costs $\bigO(d_vd_k)$ per token, independent of position. Prefill does $\bigO(Td_vd_k)$
element-wise work with a work-efficient scan (or $\bigO(Td_vd_k\log T)$ with the Hillis--Steele
doubling scan the PyTorch backend uses) in $\bigO(\log T)$ depth. Like Mamba's selective scan, and
unlike the chunkwise matrix-multiply form of GLA (\cref{ch:linear}), every operation is
element-wise per cell, so the kernels cannot use tensor-core matrix multiplications; they follow
Mamba's fused-scan design instead (keep the expanded state on chip, recompute it in the backward
pass).

\subsection{Implementation walkthrough}

\begin{implnote}[title={In the code: the KLA layer}]
The layer is \code{KLALayer} in \repofile{repos/vaisakh-shaj__kalman-linear-attention/src/kla/layers/kla_layer.py},
configured by \code{KLAConfig} in \repofile{repos/vaisakh-shaj__kalman-linear-attention/src/kla/configs.py}.
For input $[B,L,D]$ the forward pass is: \code{in_proj} to $(z,\text{gate})$, each
$[B,L,d_\text{inner}]$ with $d_\text{inner}=\text{\code{expand}}\cdot D$ (default \code{expand=2}); a
depthwise causal convolution (kernel 4, SiLU); \code{_project_sensors}, which splits
\code{sensor_proj(z)} into the value, $\log\sigma_v^2$, key and query; the scan; a learned skip
$y+\lambda_\text{skip}z$ (scalar, initialized to $-1$); SiLU output gating; and \code{out_proj}.
\begin{itemize}
  \item \textbf{Two published blocks.} The default ``Mamba-style'' block sets \code{value_rank="conv"},
  i.e.\ $\vv_t=z_t$ with no value projection, and produces $\log\sigma_v^2$ through a low-rank
  bottleneck of width $\lceil D/8\rceil$ (\code{var_rank="auto"}). The ``plain'' block used for
  the MAD synthetic tasks uses full projections (\code{value_rank="full"}, \code{var_rank="full"}).
  \item \textbf{Precision.} $\sigma_v^2=\mathrm{softplus}(\cdot)+\text{\code{obs_var_min}}$ ($10^{-4}$) and
  $\bm{\pi}_t=1/\sigma_v^2$; the variance head starts at the constant $1/d_\text{state}$ with zero
  output weights.
  \item \textbf{Keys and queries} are $\ell_2$-normalized (\code{qk_norm=True}) and have shape
  $[B,L,S]$, shared by all $M=d_\text{inner}$ channels.
  \item \textbf{Dynamics.} \code{_continuous_params} sets $a_c=-\exp(\text{\code{lambda_log}})$
  (initialized to $-1$); \code{_discretize} applies the Ornstein--Uhlenbeck formulas above
  (\code{discretization="ou"}, described in the config as ``critical for stacking layers''), with
  $\Delta=\mathrm{softplus}(\text{\code{delta}})$ initialized log-uniformly in $[10^{-3},10^{-1}]$ and process
  noise initialized to $3\abs{a_c}\Delta$.
  \item \textbf{Variance output.} With \code{return_variance=True} the layer propagates
  $\vo^{\mathrm{var}}$ through the gate (times $g^2$) and through \code{out_proj} with squared
  weights, a diagonal approximation.
\end{itemize}
\end{implnote}

\begin{implnote}[title={In the code: the reference scan}]
\code{kla_scan_torch} in \repofile{repos/vaisakh-shaj__kalman-linear-attention/src/kla/ops/kla_ops.py}
is the readable reference that every backend is validated against (together with the sequential
\code{kla_scan_reference}). Inputs are $v$, $\Lambda^v$ ($=\bm{\pi}$) of shape $[B,L,M]$, $k$,
$q$ of shape $[B,L,S]$, and $a$, $p$ of shape $[M,S]$; the state is \texttt{KLAState(lam,~eta)},
each $[B,M,S]$, initialized to the standard-normal prior (precision 1, information 0). The code
calls our information $u$ \code{eta}. \code{_sufficient_stats} forms $\phi$ and $r$
\eqref{eq:c08-kla-stats} as outer products over $(m,s)$. \Cref{lst:c08-kla-ops} shows the two scans.
Two numerical details are worth knowing. The docstring of \code{_sufficient_stats} explains that a
floor on $\phi$ is unnecessary in the linear, trace-normalized scan (a zero-information leaf is a
valid pure-prediction step) and harmful because a clamp kills the gradient to small keys; the
optional log-space combine (\code{mobius_impl="log"}) needs the floor because it takes $\log\phi$.
The generic scans live in \repofile{repos/vaisakh-shaj__kalman-linear-attention/src/kla/ops/scan.py}:
\code{associative} (PyTorch's higher-order \code{associative_scan}), \code{doubling} (vectorized
Hillis--Steele) and \code{sequential}; without gradients the reference switches to \code{doubling}
to avoid recompiling for every sequence length.
\end{implnote}

\begin{lstlisting}[caption={Excerpt from \texttt{repos/\allowbreak vaisakh-shaj\_\_kalman-linear-attention/\allowbreak src/\allowbreak kla/\allowbreak ops/\allowbreak kla\_ops.py} (function \texttt{kla\_scan\_torch}).},label={lst:c08-kla-ops}]
    if mobius_impl == "linear":
        A = ((1.0 + p_ * phi) / a2).expand_as(phi)
        C = (p_ / a2).expand_as(phi)
        D = torch.ones_like(phi)
        pA, pB, pC, pD = scan(_mobius_combine_tracenorm, (A, phi, C, D), dim=1)
        lam0_ = lam0.unsqueeze(1)
        lam = (pA * lam0_ + pB) / (pC * lam0_ + pD).clamp_min(EPS)
        var = 1.0 / lam.clamp_min(EPS)
    # ...
    lam_prev = torch.cat((lam0.unsqueeze(1), lam[:, :-1]), dim=1)
    denom = (a2 + p_ * lam_prev).clamp_min(EPS)
    alpha = a_ / denom

    pAlpha, pR = scan(_affine_combine, (alpha, r), dim=1)
    eta = pAlpha * eta0.unsqueeze(1) + pR

    mean = eta * var

    if decode_from_prior:
        mean = a_ * mean
        var = a2 * var + p_

    y = torch.einsum("blms,bls->blm", mean, q)
    y_var = torch.einsum("blms,bls->blm", var, q * q)
    return y, y_var, KLAState(lam=lam[:, -1], eta=eta[:, -1])
\end{lstlisting}

The optional \code{decode_from_prior} branch emits the one-step-ahead prediction (mean $a\cdot$mean,
variance $a^2\cdot$var$+p$) instead of the filtered posterior, i.e.\ the belief about the
\emph{next} step before its token is seen.

\begin{implnote}[title={In the code: fused GPU kernels}]
The Triton backend (\repofile{repos/vaisakh-shaj__kalman-linear-attention/src/kla/ops/triton_backend.py})
uses a single fused forward kernel when no gradient is needed and two tiled scans
(\code{tiled_mobius_scan}, \code{tiled_affine_scan}) with an exact scalar-adjoint backward otherwise.
In the fused kernel (\repofile{repos/vaisakh-shaj__kalman-linear-attention/src/kla/ops/kernels/triton/fused_kla_scan.py})
one program owns one (batch, channel) pair, streams the sequence in chunks of \code{BLOCK_L}$=64$
tokens, keeps all $S$ states of the channel in the second axis of a $[\text{\code{BLOCK_L}},S]$ tile and
carries the running $2\times2$ matrix and the information across chunks (\cref{lst:c08-kla-fused}).
Note that inside the kernel \code{q_ptr}/\code{q_t} is the process noise $p$, not the query (the
query is \code{w}), and \code{h} is the key. The ``key trick'' named in its docstring recovers
$\Lambda_{t-1}$ locally by inverting the leaf: from $\Lambda_t=(A\Lambda_{t-1}+\phi)/(C\Lambda_{t-1}+1)$
we get $\Lambda_{t-1}=(\Lambda_t-\phi)/(A-C\Lambda_t)$, so the gate needs no shift across chunk
boundaries. The CUDA backend (\code{backend="cuda"}, sources under
\repofile{repos/vaisakh-shaj__kalman-linear-attention/src/kla/ops/kernels/cuda/v3/}) uses CUB block
primitives with the same trace-normalized operator, saves only chunk-boundary checkpoints and recomputes
inside each chunk in the backward pass, avoiding a $[B,L,M,S]$ state history in memory; it supports
static dynamics with $d_\text{state}\le64$ and does not return the filter state, so stateful decoding
uses the Triton or PyTorch path. An Apple-silicon backend (\code{mps}) uses fused Metal kernels.
\end{implnote}

\begin{lstlisting}[caption={Excerpt from \texttt{repos/\allowbreak vaisakh-shaj\_\_kalman-linear-attention/\allowbreak src/\allowbreak kla/\allowbreak ops/\allowbreak kernels/\allowbreak triton/\allowbreak fused\_kla\_scan.py} (inside the chunk loop of \texttt{\_fused\_fwd\_kernel}). One non-ASCII comment line is elided.},label={lst:c08-kla-fused}]
        A = (1.0 + q_t * phi) / a2_t
        C = q_t / a2_t
        D = zero + 1.0

        sA, sB, sC, sD = tl.associative_scan(
            (A, phi, C, D), axis=0, combine_fn=_tn_combine
        )
        fA = sA * cA[None, :] + sB * cC[None, :]
        fB = sA * cB[None, :] + sB * cD[None, :]
        fC = sC * cA[None, :] + sD * cC[None, :]
        fD = sC * cB[None, :] + sD * cD[None, :]
        inv = 1.0 / tl.maximum(fA + fD, 1e-12)
        fA, fB, fC, fD = fA * inv, fB * inv, fC * inv, fD * inv

        lam = (fA * lam0[None, :] + fB) / tl.maximum(fC * lam0[None, :] + fD, 1e-12)
        # ...
        lam_prev = (lam - phi) / tl.maximum(A - C * lam, 1e-12)
        alpha = a_t / tl.maximum(a2_t + q_t * lam_prev, 1e-12)

        ga, gb = tl.associative_scan((alpha, rr), axis=0, combine_fn=_aff_combine)
        eta = ga * c_eta[None, :] + gb
\end{lstlisting}

Lines 8--13 compose the in-chunk prefix (variables \code{sA} to \code{sD}) with the matrix
carried from earlier chunks (\code{cA} to \code{cD}): the carried matrix is applied first, so it multiplies from the right, exactly the
order of $\mathcal{M}_t\cdots\mathcal{M}_1$. The repository's documentation (\repofile{repos/vaisakh-shaj__kalman-linear-attention/docs/backends.md})
states the shared accuracy test for every backend: maximum absolute error at most
$2\times10^{-5}+2\times10^{-4}\cdot\max\abs{\text{reference}}$ against the sequential reference, for
outputs and input gradients.

\subsection{Reference implementation}

\Cref{lst:c08-kla-ours} is a minimal version of the whole mechanism: the two combine rules, a
Hillis--Steele scan, the parallel layer core and the decode step. We ran it through a small NumPy
stand-in for the few \code{torch} calls it uses: the parallel output and the token-by-token decode
agreed to $10^{-15}$, as did the final states.

\begin{lstlisting}[caption={Reference implementation (ours): the KLA core as two scans, plus the decode step.},label={lst:c08-kla-ours}]
import torch

def mobius(l, r):                       # compose 2x2 maps (r after l), trace-normalised
    A = r[0] * l[0] + r[1] * l[2]; B = r[0] * l[1] + r[1] * l[3]
    C = r[2] * l[0] + r[3] * l[2]; D = r[2] * l[1] + r[3] * l[3]
    s = A + D
    return A / s, B / s, C / s, D / s

def affine(l, r):                       # compose maps z -> a z + b (r after l)
    return r[0] * l[0], r[0] * l[1] + r[1]

def scan(combine, xs):                  # inclusive Hillis-Steele scan over dim 1
    xs, n, off = list(xs), xs[0].shape[1], 1
    while off < n:
        new = combine([x[:, :-off] for x in xs], [x[:, off:] for x in xs])
        xs = [torch.cat([x[:, :off], y], 1) for x, y in zip(xs, new)]
        off *= 2
    return xs

def kla(v, pv, k, q, a, p, lam0, u0):
    """v, pv: [B,T,M] values and value precisions; k, q: [B,T,S] keys, queries;
    a, p: [M,S] decay and process noise; lam0, u0: [B,M,S] prior precision, information."""
    phi = pv[..., None] * (k * k)[:, :, None]          # information added   [B,T,M,S]
    r = (pv * v)[..., None] * k[:, :, None]             # information-vector input
    a2 = a * a
    A = (1 + p * phi) / a2                              # leaf [[A, phi], [C, 1]]
    C = (p / a2).expand_as(phi)
    pA, pB, pC, pD = scan(mobius, (A, phi, C, torch.ones_like(phi)))
    lam = (pA * lam0[:, None] + pB) / (pC * lam0[:, None] + pD)      # precision
    lam_prev = torch.cat([lam0[:, None], lam[:, :-1]], 1)
    alpha = a / (a2 + p * lam_prev)                     # uncertainty-dependent gate
    gA, gB = scan(affine, (alpha, r))
    u = gA * u0[:, None] + gB                           # information vector
    y = torch.einsum('btms,bts->btm', u / lam, q)       # read posterior mean
    y_var = torch.einsum('btms,bts->btm', 1 / lam, q * q)
    return y, y_var, (lam[:, -1], u[:, -1])

def kla_step(v, pv, k, q, a, p, lam, u):              # one decode step: O(MS)
    den = a * a + p * lam
    lam = lam / den + pv[..., None] * (k * k)[:, None]
    u = (a / den) * u + (pv * v)[..., None] * k[:, None]
    return torch.einsum('bms,bs->bm', u / lam, q), (lam, u)
\end{lstlisting}

\subsection{Reported results}

\begin{results}
The repository contains no benchmark tables, so we report only what its files state and refer to
\citet{shaj2026kalman} for numbers.
\begin{itemize}
  \item \repofile{repos/vaisakh-shaj__kalman-linear-attention/readme.md} positions KLA against
  softmax attention and SSMs/GLA: ``fractional linear (M\"obius)'' expressivity versus linear for
  SSMs/GLA, $\bigO(T)$ training and $\bigO(1)$ per-token inference, explicit sequence uncertainty,
  and parallel training. It describes KLA as a drop-in replacement for other linear attention layers.
  \item \repofile{repos/vaisakh-shaj__kalman-linear-attention/docs/usage.md}: at $d_\text{model}=512$
  the Mamba-style block has about 1.79M parameters against 3.76M for the plain block; both use the
  same Kalman scan. The paper uses the plain block for MAD and the Mamba-style block for FineWeb-Edu
  pretraining.
  \item \repofile{repos/vaisakh-shaj__kalman-linear-attention/docs/backends.md}: the CUDA backend
  was the fastest in the paper's NVIDIA GPU training benchmarks, ``substantially faster than
  Triton'', with a much smaller scan-memory footprint.
  \item The survey \citep{zhoubian2026memory} describes Kalman Linear Attention as using filtering
  ``to modulate updates by estimation quality''.
\end{itemize}
Experiment code is marked ``Coming Soon'' in the README; see the paper for language-modelling and
state-tracking results.
\end{results}

\subsection{Discussion}

\textbf{Strengths.} KLA turns a principled estimator into a layer without giving up parallel
training: the information form makes writes additive, and the leftover nonlinearity (the precision
predict step) is linear-fractional and therefore scannable. The write gate is not a free parameter
but follows from accumulated evidence, and the layer returns a variance estimate at no extra
cost. Process noise gives a clean, interpretable account of forgetting.

\textbf{Limitations.} The filter is exact only for the factorized model \eqref{eq:c08-kla-model}:
cells do not share uncertainty, so KLA has no analogue of the key--key coupling
$\vk_t\vk_t\T$ through which the delta rule (\cref{ch:delta}) erases the old value stored at a
key. Its transition is diagonal and time-invariant ($a$, $p$, $\Delta$ are static), so its
expressivity comes from the precision-dependent gate rather than from richer transitions such as
those of \cref{ch:expressive}. The state is twice as large as GLA's, and the element-wise scan
cannot use tensor cores.

\textbf{Relations.} With the precision frozen at its steady state, KLA reduces to GLA with a
constant gate (\cref{ch:linear}); with $p=0$ it reduces to a non-forgetting accumulator. GKA,
next, keeps the full key covariance that KLA factorizes away.

% ===========================================================================================
\section{Gated KalmaNet}\label{sec:c08-gka}

\subsection{Motivation}

Linear SSMs and linear-attention layers update their state from the current token alone: a decay,
an outer product, at most one delta-rule correction. The resulting state is a lossy, fading
summary of the past, and recall-heavy tasks expose the loss. \Citet{peng2026gka} ask instead for
the state that is \emph{optimal given the entire past}: at every token, the solution of a ridge
regression over all previous key--value pairs, weighted by a gated, fading memory. The README of the
toolkit that ships the official code puts it this way: GKA is ``Inspired by Kalman filtering, it
computes the state at each time-step based on the entire past rather than the instantaneous update
rules used by other SSMs'' (\repofile{repos/awslabs__hybrid-model-factory/README.md}). The authors'
abstract describes the reduction: grounding the layer in the Kalman filter and maintaining the full
error covariance, a steady-state assumption that enables parallelization turns the problem into
online ridge regression with constant memory and linear compute. Doing this in practice requires
solving two engineering problems that the textbook recursions of \cref{sec:c08-rls} do not: they
are numerically fragile in bfloat16 and inherently sequential. GKA's answers are an input-dependent
(adaptive) regularizer that bounds the condition number, Chebyshev iteration instead of RLS or
conjugate gradient, and a chunked kernel.

\begin{taxonomy}
In the taxonomy of \citet{zhoubian2026memory}, GKA is \textbf{implicit} memory (its state is part
of the layer's forward computation), with \textbf{online} update dynamics and \textbf{long-term}
persistence within a sequence (survey Table~I groups it with Kalman Linear Attention as
``Probabilistic Recurrent State Estimation''). Its \textbf{update rule} is a \emph{state-transition
update} (survey Table~II): the survey says Kalman Linear Attention and Gated KalmaNet ``use
filtering or online regression to modulate updates by estimation quality''. Unlike every other
layer in \cref{ch:linear,ch:delta}, the ``transition'' is not a fixed function of the current token
but the solution of an optimization problem over the whole gated past.
\end{taxonomy}

\subsection{Formulation: test-time ridge regression}\label{sec:c08-gka-form}

We use GKA's notation mapped to ours. Per head, keys and queries are $\ell_2$-normalized; a
\emph{beta gate} $\beta_t\in(0,1)$ multiplies both $\vk_t$ and $\vv_t$; a forgetting gate
$\alpha_t=e^{g_t}\in(0,1]$ decays the past (the code's \code{g}, computed Mamba-2 style); the
regularization ratio $r$ (the code's \code{ridge_strength}, default $0.02$) sets the ridge strength;
and a mixing gate $\nu_t\in(0,1)$ (the code calls it the ``alpha connection'', \code{alpha_proj})
blends two read-outs. The code's module docstring
(\repofile{repos/awslabs__hybrid-model-factory/training/src/hmf/model/hybrid_zoo/layers/gated_kalmanet/gka.py})
states the objective: at each time $t$
\begin{equation}\label{eq:c08-gka-obj}
  \mS_t=\argmin_{\mS}\ \lambda_t\norm{\mS}_F^2+\sum_{i\le t}\eta_{t,i}\norm{\mS\vk_i-\vv_i}^2,
  \qquad \eta_{t,i}=\beta_i^2\prod_{j=i+1}^{t}\alpha_j ,
\end{equation}
where $\vk_i,\vv_i$ are the normalized key and the value before the beta gate (applying $\beta_i$ to
both $\vk_i$ and $\vv_i$ multiplies the squared error by $\beta_i^2$). By \cref{sec:c08-ewridge} the
solution is
\begin{equation}\label{eq:c08-gka-sol}
  \mS_t=\mU_t(\mH_t+\lambda_t\mI)^{-1},\qquad
  \mH_t=\alpha_t\mH_{t-1}+\beta_t^2\vk_t\vk_t\T,\qquad
  \mU_t=\alpha_t\mU_{t-1}+\beta_t^2\vv_t\vk_t\T .
\end{equation}

\paragraph{The Kalman reading.} One way to see the connection (ours; the paper's derivation via a
steady-state Kalman filter may differ in detail): the static-state Kalman filter of
\cref{sec:c08-rls} with prior $\mathcal{N}(\bm{0},\lambda_t^{-1}\mI)$ on each row of $\mS$, in which
the observation of pair $i$ is treated at time $t$ as having noise variance $1/\eta_{t,i}$, has
exactly \eqref{eq:c08-gka-sol} as its posterior mean (its negative log-posterior is half the
objective in \eqref{eq:c08-gka-obj}). Fading memory is then ``old observations
become noisier'', and the Kalman gain uses the full posterior covariance
$(\mH_t+\lambda_t\mI)^{-1}$ instead of the delta rule's $\beta_t\mI$.

\paragraph{The read-out trick.} Forming $\mS_t$ would require a $d_k\times d_k$ inverse per token.
But the layer only needs $\mS_t\vq_t$, and
\begin{equation}\label{eq:c08-gka-read}
  \mS_t\vq_t=\mU_t\tilde\vq_t,\qquad \tilde\vq_t=(\mH_t+\lambda_t\mI)^{-1}\vq_t .
\end{equation}
So GKA (i) solves one $d_k$-dimensional linear system per token for the \emph{whitened query}
$\tilde\vq_t$, then (ii) reads the gated-linear-attention memory $\mU_t$ with it. Step (ii) is plain
GLA. The full output is
\begin{equation}\label{eq:c08-gka-out}
  \hat\vq_t=\nu_t\tilde\vq_t+(1-\nu_t)\vq_t,\qquad \vo_t=\tfrac{1}{\sqrt{d_k}}\,\mU_t\hat\vq_t ,
\end{equation}
followed by a gated RMSNorm and the output projection. With $\nu_t=1$ the layer returns the ridge
read; with $\nu_t=0$ it is exactly GLA. The learned per-token, per-head $\nu_t$ lets the model choose.

\begin{example}[Why whitening the query helps]\label{ex:c08-whiten}
Let $d_k=2$, $d_v=1$, no decay, $\beta=1$, and store two nearly parallel keys
$\vk_1=(1,0)$, $\vk_2=(\cos30^\circ,\sin30^\circ)\approx(0.866,0.5)$ with values $v_1=1$, $v_2=-1$.
Query with $\vq=\vk_1$, hoping to retrieve $v_1=1$.
\begin{itemize}
  \item Linear attention: $\mU\vk_1=v_1(\vk_1\T\vk_1)+v_2(\vk_2\T\vk_1)=1-0.866=0.134$. The
  similar key interferes.
  \item Delta rule with $\beta=1$, writing $\vk_1$ then $\vk_2$: the second write fixes
  $\mS\vk_2=-1$ exactly but drags $\mS\vk_1$ to $-0.616$.
  \item Ridge with $\lambda\to0$: $\mH=\vk_1\vk_1\T+\vk_2\vk_2\T$ is invertible, and $\mS=\mU\mH^{-1}$
  satisfies $\mS\vk_1=1$ and $\mS\vk_2=-1$ exactly. The whitened query
  $\tilde\vq=\mH^{-1}\vk_1$ is the dual vector that has inner product 1 with $\vk_1$ and 0 with
  $\vk_2$, so it cancels the interference.
  \item GKA's adaptive ridge, $\lambda=0.02\norm{\mH}_F\approx0.037$: the read is $0.78$. The
  eigenvalues of $\mH$ are $1.87$ and $0.134$; the small one (the direction that separates the two
  keys) is shrunk by the factor $0.134/(0.134+0.037)$. Regularization trades exact recall of
  near-duplicate keys for stability.
\end{itemize}
\end{example}

\subsection{Adaptive regularization and the Chebyshev solver}\label{sec:c08-gka-cheb}

GKA sets the ridge strength proportional to the size of the key covariance,
\begin{equation}\label{eq:c08-gka-lambda}
  \lambda_t=r\,\norm{\mH_t}_F .
\end{equation}
\begin{proposition}[Bounded condition number]\label{prop:c08-kappa}
With \eqref{eq:c08-gka-lambda}, the eigenvalues of $\mH_t+\lambda_t\mI$ lie in
$[\lambda_t,\ \lambda_t+\norm{\mH_t}_F]$, so $\kappa(\mH_t+\lambda_t\mI)\le1+1/r$ for every input.
\end{proposition}
\begin{proof}
$\mH_t$ is a nonnegative combination of outer products, so it is symmetric positive semidefinite
with eigenvalues $0\le\mu_i$. Its largest eigenvalue satisfies $\mu_{\max}\le(\sum_i\mu_i^2)^{1/2}
=\norm{\mH_t}_F$. Adding $\lambda_t\mI$ shifts every eigenvalue by $\lambda_t$, so the spectrum lies
in $[\lambda_t,\lambda_t+\norm{\mH_t}_F]$ and the ratio of the ends is
$1+\norm{\mH_t}_F/\lambda_t=1+1/r$.
\end{proof}
The bound is tight when $\mH_t$ has rank one; in our check over 200 random key sets the largest
condition number observed was exactly 51.0 for $r=0.02$. Three things follow.
\begin{enumerate}
  \item \textbf{Stability for any input.} However the gates scale the memory (a long effective
  window makes $\mH_t$ large, a short one makes it small), the linear system has $\kappa\le51$, a
  regime in which bfloat16 arithmetic behaves.
  \item \textbf{Scale invariance.} Multiplying all weights $\eta_{t,i}$ by a common factor $c$
  multiplies $\mH_t$, $\mU_t$ and $\lambda_t$ by $c$ and leaves $\mS_t$ unchanged, so the balance
  between data and regularizer does not depend on the overall magnitude of the gates.
  \item \textbf{A known spectral interval for free.} Chebyshev iteration needs
  $[\mu_{\min},\mu_{\max}]$; the proposition supplies $[\lambda_t,\lambda_t+\norm{\mH_t}_F]$ at the
  cost of one Frobenius norm.
\end{enumerate}
Plugging this interval into \cref{sec:c08-solvers} gives step size
$\tau_t=2/(\mu_{\min}+\mu_{\max})=2/(2\lambda_t+F_t)$ and
$\rho_t=(\mu_{\max}-\mu_{\min})/(\mu_{\max}+\mu_{\min})=F_t/(2\lambda_t+F_t)=1/(1+2r)$, with
$F_t=\norm{\mH_t}_F$; these are literally the variables \code{stepsize} and \code{rho_sq_4}
($=\rho^2/4$) of the kernels in \cref{lst:c08-gka-cheb}. With $\kappa\le51$,
$\sigma=(\sqrt{51}-1)/(\sqrt{51}+1)\approx0.754$, and the worst-case relative error after $n$
iterations is $2\sigma^n$: about $0.49$ after 5 iterations, $0.12$ after 10 and $4.2\times10^{-4}$
after 30, the default. The default thus drives the worst-case solver error below bfloat16's unit
roundoff ($2^{-8}\approx3.9\times10^{-3}$); smaller \code{num_iter} values trade accuracy for
speed, and the worst case is pessimistic (in \cref{ex:c08-race} the error after 10 iterations was
$7\times10^{-2}$ and after 30 iterations $2.4\times10^{-4}$).

Why Chebyshev and not the textbook recursions? RLS cannot be used with a fixed, non-decaying
regularizer and gating (\cref{sec:c08-ewridge}), and its covariance update subtracts nearly equal
quantities, which destroys positive definiteness in low precision. Conjugate gradient needs inner
products whose rounding errors change its coefficients. The authors state that Chebyshev iteration
is more stable than conventional iterative solvers in low precision, which is consistent with its
data-independent coefficients; our tiny float16 simulation (\cref{ex:c08-race}) was too small to
show a difference, so see the paper for the evidence.

\subsection{Chunkwise computation}\label{sec:c08-gka-chunk}

During training every token needs products $\mH_t\vx$ for its own matrix $\mH_t$, in each of the
$n_\text{iter}$ iterations. Materializing $\mH_t$ for every token would cost $Td_k^2$ memory per
head. GKA uses the chunk decomposition of \cref{ch:linear}. Split the sequence into chunks of $C$
tokens, let $\mH_{[i]}$ be the state at the start of chunk $i$, and let $G_t=\sum_{j=\text{start}}^{t}
g_j$ be the chunk-local cumulative log-decay. Then for $t$ in chunk $i$ (keys already include
$\beta$)
\begin{equation}\label{eq:c08-gka-chunkH}
  \mH_t = e^{G_t}\mH_{[i]}+\sum_{j\le t,\ j\in[i]}e^{G_t-G_j}\vk_j\vk_j\T
  \quad\Longrightarrow\quad
  \mH_t\vx = e^{G_t}\mH_{[i]}\vx+\sum_{j\le t}e^{G_t-G_j}(\vk_j\T\vx)\,\vk_j .
\end{equation}
Stack the $C$ current iterates as rows of $\mX\in\R^{C\times d_k}$ and the keys as rows of
$\mK_{[i]}$. With the decay mask $D_{tj}=e^{G_t-G_j}$ for $j\le t$ and $0$ otherwise, all $C$
products are
\begin{equation}\label{eq:c08-gka-chunkHX}
  \bigl[\mH_t\vx_t\bigr]_{t\in[i]} = \diag(e^{G})\,\mX\mH_{[i]}+\bigl((\mX\mK_{[i]}\T)\had D\bigr)\mK_{[i]} ,
\end{equation}
two matrix multiplications of sizes $C\times d_k\times d_k$ and $C\times C\times d_k$ (using the
symmetry of $\mH_{[i]}$). The Frobenius norms need the same treatment. Expanding the square of
\eqref{eq:c08-gka-chunkH} with $\inner{\mH}{\vk\vk\T}_F=\vk\T\mH\vk$ and
$\inner{\vk\vk\T}{\vk'\vk'^{\top}}_F=(\vk\T\vk')^2$,
\begin{equation}\label{eq:c08-gka-fro}
  \norm{\mH_t}_F^2=e^{2G_t}\norm{\mH_{[i]}}_F^2+2e^{G_t}\sum_{j\le t}D_{tj}\,\vk_j\T\mH_{[i]}\vk_j
  +\sum_{j,l\le t}D_{tj}D_{tl}(\vk_j\T\vk_l)^2 .
\end{equation}
In our NumPy check both \eqref{eq:c08-gka-chunkHX} and \eqref{eq:c08-gka-fro} matched the directly
formed $\mH_t$ to $10^{-15}$. The chunk-start states $\mH_{[i]}$ are themselves a GLA chunk
recurrence with values equal to keys, and the final read \eqref{eq:c08-gka-out} is a chunkwise GLA
call, so GKA reuses the GLA machinery of \cref{ch:linear} twice and adds the in-chunk solver.

\paragraph{Backward pass.} Treat $\tilde\vq=\mA^{-1}\vq$ with $\mA=\mH+\lambda\mI$ symmetric. For a
loss $\loss$ with incoming gradient $\bar{\vx}=\partial\loss/\partial\tilde\vq$, implicit
differentiation gives
\begin{equation}\label{eq:c08-gka-bwd}
  \frac{\partial\loss}{\partial\vq}=\mA^{-1}\bar{\vx},\qquad
  \frac{\partial\loss}{\partial\mA}=-\Bigl(\frac{\partial\loss}{\partial\vq}\Bigr)\tilde\vq\T .
\end{equation}
The first is the same linear solve applied to the gradient, so the backward pass reruns the
Chebyshev solver; the second is an outer product that flows into the keys through
\eqref{eq:c08-gka-sol}. Because $\lambda=r\norm{\mH}_F$ also depends on $\mH$, there is an extra
term: $\partial\loss/\partial\lambda=\tr(\partial\loss/\partial\mA)=-\bar{\vq}\T\tilde\vq$ (with
$\bar{\vq}=\partial\loss/\partial\vq$) and $\partial\norm{\mH}_F/\partial\mH=\mH/\norm{\mH}_F$, so
$\mH$ receives the additional gradient $-r(\bar{\vq}\T\tilde\vq)\,\mH/\norm{\mH}_F$; the flag
\code{bp_lambda} turns this term on.

\subsection{Algorithms and complexity}

\begin{algorithm}[t]
\caption{GKA prefill / training for one head (chunkwise)}\label{alg:c08-gka-prefill}
\begin{algorithmic}[1]
\Require $\vq_t,\vk_t,\vv_t$, log-decay $g_t$, gates $\beta_t,\nu_t$; ratio $r$; chunk size $C$; $n_\text{iter}$
\State $\vq_t\gets\vq_t/\norm{\vq_t}$, \ $\vk_t\gets\beta_t\vk_t/\norm{\vk_t}$, \ $\vv_t\gets\beta_t\vv_t$
\State $G\gets$ chunk-local cumulative sum of $g$
\State $\mH_{[i]}$ for all chunks $\gets$ GLA chunk-state recurrence with values $=$ keys
\For{each chunk $i$ \textbf{in parallel}}
  \State $F_t\gets\norm{\mH_t}_F$ by \eqref{eq:c08-gka-fro}; \ $\lambda_t\gets rF_t$; \ $\tau_t\gets2/(2\lambda_t+F_t)$; \ $\rho_t\gets F_t/(2\lambda_t+F_t)$
  \State $\mX_\text{prev}\gets\bm{0}$; \ $\mX\gets\diag(\tau)\mQ_{[i]}$; \ $\omega\gets2$
  \For{$n=1..n_\text{iter}$}
    \State $\omega\gets1/(1-\rho^2\omega/4)$ \Comment{per token}
    \State $\bm{R}\gets\diag(e^{G})\mX\mH_{[i]}+((\mX\mK_{[i]}\T)\had D)\mK_{[i]}+\diag(\lambda)\mX-\mQ_{[i]}$
    \State $(\mX,\mX_\text{prev})\gets\bigl(\omega\had(\mX-\diag(\tau)\bm{R})+(1-\omega)\had\mX_\text{prev},\ \mX\bigr)$
  \EndFor
  \State $\hat\mQ_{[i]}\gets\diag(\nu)\mX+\diag(1-\nu)\mQ_{[i]}$
\EndFor
\State $\mO\gets$ chunkwise GLA$(\hat\mQ,\mK,\mV,g)$ with scale $1/\sqrt{d_k}$ \Comment{\cref{ch:linear}}
\end{algorithmic}
\end{algorithm}

\begin{algorithm}[t]
\caption{GKA decode (one token, one head)}\label{alg:c08-gka-decode}
\begin{algorithmic}[1]
\Require state $(\mH,\mU)$; new $\vq,\vk,\vv,g,\beta,\nu$ (normalized and $\beta$-scaled as above)
\State $\mH\gets e^{g}\mH+\vk\vk\T$; \ $F\gets\norm{\mH}_F$; \ $\lambda\gets rF$
\State $\tilde\vq\gets$ \Call{Chebyshev}{$\mH+\lambda\mI$, $\vq$, interval $[\lambda,\lambda+F]$, $n_\text{iter}$}
\State $\mU\gets e^{g}\mU+\vv\vk\T$
\State \Return $\vo=\mU(\nu\tilde\vq+(1-\nu)\vq)/\sqrt{d_k}$ and the new state
\end{algorithmic}
\end{algorithm}

\Cref{alg:c08-gka-prefill,alg:c08-gka-decode} summarize GKA. Per head the recurrent state is
$\mH_t$ and $\mU_t$, $2d_k^2$ numbers for $d_v=d_k$, twice a GLA or Gated DeltaNet head. Each
Chebyshev iteration on a chunk costs $\bigO(Cd_k^2+C^2d_k)$, so prefill costs
$\bigO\bigl((n_\text{iter}+1)\,T(d_k^2+Cd_k)\bigr)$ per head, roughly $n_\text{iter}+1$ times GLA,
with all iterations running on chip inside one kernel per (chunk, head). Decoding costs
$\bigO(n_\text{iter}d_k^2)$ per token per head, independent of position. The iteration count is a
\emph{test-time compute knob}: it can be lowered at serving time without retraining, trading
accuracy of the solve for speed.

\subsection{Implementation walkthrough}

\begin{implnote}[title={In the code: the GKA layer}]
GKA ships inside the AWS Hybrid Model Factory toolkit. The layer is \code{GatedKalmaNet} in
\repofile{repos/awslabs__hybrid-model-factory/training/src/hmf/model/hybrid_zoo/layers/gated_kalmanet/gka.py},
wrapped by \code{GatedKalmaNetLayer} in \code{modules.py} of the same folder, and configured by
\code{GKAConfig} in \code{gka_config.py}. Its hyperparameters are chosen so that the layer can be
initialized from an existing attention layer: \code{_construct_from_transformer} copies the head
dimension and the (grouped-query) head counts of the base Transformer. In \code{forward}:
\begin{itemize}
  \item $q,k,v$ projections, each followed by a \code{ShortConvolution} (kernel 4, SiLU, from
  \fla); key/value heads are repeated (or expanded by an optional low-rank projection) to the
  query head count;
  \item $\ell_2$ normalization of $q$ and $k$ (epsilon $10^{-6}$), then the beta gate, the
  alpha connection $\nu$ and the forgetting gate (\cref{lst:c08-gka-gates});
  \item \code{gka_chebyshev_gla} (training/prefill) or \code{torch_decoding_one_step} (decoding),
  then \code{FusedRMSNormGated} with a gate from \code{g_proj}, and \code{o_proj}.
\end{itemize}
The cache stores $(\mH_t,\mU_t)$ as \texttt{(h\_kk,~h\_kv)} plus the three convolution states. Defaults
(\code{GKAConfig}): \code{num_iter=30}, \code{ridge_strength=0.02} (our $r$), \code{chunk_size=64},
\code{gla_rescale=True} (scale $1/\sqrt{d_k}$), \code{use_forgetting_gate_kk=True} (the same decay
$g$ for $\mH$ as for $\mU$), \code{bp_lambda=True}.
\end{implnote}

\begin{lstlisting}[caption={Excerpt from \texttt{repos/\allowbreak awslabs\_\_hybrid-model-factory/\allowbreak training/\allowbreak src/\allowbreak hmf/\allowbreak model/\allowbreak hybrid\_zoo/\allowbreak layers/\allowbreak gated\_kalmanet/\allowbreak gka.py} (method \texttt{GatedKalmaNet.forward}).},label={lst:c08-gka-gates}]
        # Apply beta gating if enabled
        if self.use_beta_gate:
            beta = self.b_proj(hidden_states).sigmoid()  # [B_sp, Lp, H]
            eps = 1e-6
            k = ((beta[..., None] + eps) * k).to(k.dtype)  # [B_sp, Lp, H, D]
            v = ((beta[..., None] + eps) * v).to(v.dtype)  # [B_sp, Lp, H, D]

        # Compute alpha connection parameter
        if self.use_alpha_connection:
            alpha = self.alpha_proj(hidden_states).sigmoid()  # [B_sp, Lp, H]
        else:
            alpha = torch.ones_like(q[..., 0])  # [B_sp, Lp, H]

        # Compute forgetting gate
        if self.use_forgetting_gate:
            g = -self.A_log.float().exp() * F.softplus(
                self.a_proj(hidden_states).float() + self.dt_bias
            )  # [B_sp, Lp, H]
        else:
            g = None

        gk = g if self.use_forgetting_gate_kk else None
\end{lstlisting}

The decay is Mamba-2's parameterization, $g_t=-e^{A_{\log}}\,\mathrm{softplus}(\vw_a\T\vx_t+b_{\Delta})$,
with $A=e^{A_{\log}}$ initialized uniformly in $(0,16)$ and $b_\Delta$ initialized so that the step
lies in $[10^{-3},10^{-1}]$. The training-time entry point then does two things
(\cref{lst:c08-gka-solve}): solve for the whitened queries and mix them with the raw queries, and
hand everything to \fla's \code{chunk_simple_gla}.

\begin{lstlisting}[caption={Excerpt from \texttt{repos/\allowbreak awslabs\_\_hybrid-model-factory/\allowbreak training/\allowbreak src/\allowbreak hmf/\allowbreak model/\allowbreak hybrid\_zoo/\allowbreak layers/\allowbreak gated\_kalmanet/\allowbreak ops/\allowbreak chebyshev/\allowbreak gka\_chebyshev\_solve.py} (function \texttt{gka\_chebyshev\_gla}). Two non-ASCII comments are elided.},label={lst:c08-gka-solve}]
    if num_iter > 0:
        # ...
        q_, kk_final = latent_chebyshev(
            k, q, ridge_strength, num_iter, solver_type, bp_lambda, gk, chunk_size
        )

        q = q + alpha[..., None] * (q_ - q)  # ...
    else:
        kk_final = None

    o, kv_final = chunk_simple_gla(
        q=q, k=k, v=v, scale=gla_scale, g=g, output_final_state=True
    )

    return o, kk_final, kv_final
\end{lstlisting}

\begin{implnote}[title={In the code: the chunked Chebyshev kernel}]
\code{latent_chebyshev} calls the autograd function \code{ChebyshevIteration} in
\repofile{repos/awslabs__hybrid-model-factory/training/src/hmf/model/hybrid_zoo/layers/gated_kalmanet/ops/chebyshev/chebyshev_iteration.py}.
Its forward computes the chunk-local cumulative decay with \fla's \code{chunk_local_cumsum}, the
chunk-start states $\mH_{[i]}$ with \fla's \code{chunk_fwd_h} called with \code{v=k}, and launches
\code{chebyshev_iteration_fro_gating_forward_kernel} on a grid of (number of chunks, batch$\times$heads).
Each program loads one chunk's $q$ and $k$ ($[C,D]$) and its $\mH_{[i]}$ ($[D,D]$), computes the
Frobenius norms by \eqref{eq:c08-gka-fro} (stored for the backward pass), and runs all iterations
(\cref{lst:c08-gka-cheb}). The two \code{tl.dot} lines are \eqref{eq:c08-gka-chunkHX}: \code{mask}
holds $D_{tj}=\exp(G_t-G_j)$ and \code{exp_gkk} holds $e^{G_t}$. The update
\texttt{x~=~omega~*~x~-~grad\_and\_x\_prev} expands to
$\omega(\vx-\tau\bm{R})+(1-\omega)\vx_\text{prev}$, i.e.\ \eqref{eq:c08-cheb-iter}. The backward
(\code{ChebyshevIteration.backward}) reruns the same forward kernel on the incoming gradient with the
stored norms (\code{load_fro=1}), which is the first identity of \eqref{eq:c08-gka-bwd}, then turns
$-\bar{\vq}\tilde\vq\T$ into key gradients with \fla's \code{chunk_bwd_dh} and the custom
\code{chunk_bwd_dk_gating}; with \code{bp_lambda} the latter adds the term from
$\lambda=r\norm{\mH}_F$ (in \code{chunk_dk_backward.py} it appears as weights
\texttt{ridge\_ratio~*~sum(dq~*~x\_out)~/~fro}).
\end{implnote}

\begin{lstlisting}[caption={Excerpt from \texttt{repos/\allowbreak awslabs\_\_hybrid-model-factory/\allowbreak training/\allowbreak src/\allowbreak hmf/\allowbreak model/\allowbreak hybrid\_zoo/\allowbreak layers/\allowbreak gated\_kalmanet/\allowbreak ops/\allowbreak chebyshev/\allowbreak chebyshev\_iteration.py} (end of \texttt{chebyshev\_iteration\_fro\_gating\_forward\_kernel}).},label={lst:c08-gka-cheb}]
    ridge_strength = ridge_ratio * fro

    stepsize = 2 / (2 * ridge_strength + fro)
    rho_sq_4 = fro / (2 * ridge_strength + fro) / 2.0
    rho_sq_4 = rho_sq_4 * rho_sq_4

    x_prev = tl.zeros_like(q).to(tl.float32)
    x = stepsize * q.to(tl.float32)
    omega = tl.full((chunk_size, 1), 2.0, dtype=tl.float32)

    for _ in range(num_iter):
        omega = 1.0 / (1.0 - rho_sq_4 * omega)

        grad_and_x_prev = (
            stepsize
            * omega
            * (
                tl.dot((x * exp_gkk).to(h_kk.dtype), h_kk)
                + tl.dot((tl.dot(x.to(k.dtype), tl.trans(k)) * mask).to(k.dtype), k)
                + ridge_strength * x
                - q
            )
            + (omega - 1) * x_prev
        )
        x_prev = x

        x = omega * x - grad_and_x_prev
\end{lstlisting}

\begin{implnote}[title={In the code: decoding, sequence parallelism, serving, priming}]
\begin{itemize}
  \item \textbf{Decoding.} \code{torch_decoding_one_step} calls
  \code{chebyshev_iteration_forward_triton_inference}, whose kernel (grid batch$\times$heads)
  updates \code{h_kk} $\gets e^{g}$\code{h_kk}$+\vk\vk\T$, computes the norm and runs the same
  iteration with $D\times D$ matrix--vector products, then \fla's \code{fused_recurrent_simple_gla}
  updates $\mU$ and produces the output (\cref{alg:c08-gka-decode}).
  \item \textbf{Sequence parallelism.} Because $\mH_t$ is a gated sum, chunk states can be passed
  between GPUs like any linear-attention state: \code{state_passing_gka_p2p} in
  \code{sequence_parallel/gka_sp_utils.py} passes both \code{h_kk} and \code{h_kv}, after which
  each GPU recomputes its part with the received initial states.
  \item \textbf{Serving.} The vLLM plugin has its own grouped-query kernels under
  \repofile{repos/awslabs__hybrid-model-factory/vllm-inference/src/primed_vllm/gka/}; the iteration
  count can be overridden at serving time through vLLM's \code{--hf-overrides} flag with a
  \code{gka_overrides} entry such as \code{num_iter} $=10$
  (\repofile{repos/awslabs__hybrid-model-factory/docs/Inference.md}).
  \item \textbf{Priming.} The stage-0 config
  \repofile{repos/awslabs__hybrid-model-factory/training/examples/priming/stage0/qwen3_8b_gka.yaml}
  (pattern \texttt{*DA-GKA*-*DA-GKA*-...}) builds a fused Qwen3-8B hybrid in which every other
  layer is a ``fused GKA + Attention'' layer: the new GKA path runs in parallel with the original
  attention so that stage~1 can distill attention into GKA (\repofile{repos/awslabs__hybrid-model-factory/docs/Priming.md}),
  with the defaults above (\texttt{num\_iter:~30}, \texttt{ridge\_strength:~0.02}, all gates on). Priming
  itself, the conversion of a pretrained Transformer into a hybrid, is the subject of \cref{ch:hybrid}
  \citep{chattopadhyay2026priming}.
\end{itemize}
\end{implnote}

\subsection{Reference implementation}

\Cref{lst:c08-gka-ours} gives a recurrent GKA head (the semantics of \cref{alg:c08-gka-decode}
applied token by token) and the chunked product \eqref{eq:c08-gka-chunkHX}. Run through our NumPy
stand-in for \code{torch}, the recurrent head with 80 iterations matched the closed form
$\mU_t(\mH_t+\lambda_t\mI)^{-1}\vq_t/\sqrt{d_k}$ to $2\times10^{-10}$ (with $\nu_t=1$), and
\code{chunk_Hx} matched the directly formed $\mH_t\vx_t$ to $10^{-15}$.

\begin{lstlisting}[caption={Reference implementation (ours): a recurrent GKA head and the chunked product.},label={lst:c08-gka-ours}]
import torch

def chebyshev_solve(H, q, lam, fro, n_iter):
    """Solve (H + lam I) x = q; the eigenvalues lie in [lam, lam + fro]."""
    step = 2.0 / (2 * lam + fro)                        # 2 / (mu_min + mu_max)
    rho_sq_4 = (fro / (2 * lam + fro) / 2) ** 2         # (rho / 2)^2
    x_prev, x, w = torch.zeros_like(q), step * q, 2.0
    for _ in range(n_iter):
        w = 1.0 / (1.0 - rho_sq_4 * w)                  # Chebyshev weight omega_n
        r = H @ x + lam * x - q                         # residual = gradient
        x, x_prev = w * (x - step * r) + (1 - w) * x_prev, x
    return x

def gka_recurrent(q, k, v, g, beta, nu, ratio=0.02, n_iter=30):
    """One head. q, k, v: [T,D]; g: [T] log-decay (<= 0); beta, nu: [T] in (0,1)."""
    T, D = q.shape
    q = q / q.norm(dim=-1, keepdim=True)
    k = k / k.norm(dim=-1, keepdim=True)
    k, v = beta[:, None] * k, beta[:, None] * v         # sample weight beta^2
    H, U, out = torch.zeros(D, D), torch.zeros(v.shape[1], D), []
    for t in range(T):
        a = torch.exp(g[t])
        H = a * H + torch.outer(k[t], k[t])             # fading key covariance H_t
        U = a * U + torch.outer(v[t], k[t])             # fading value-key memory U_t
        fro = torch.linalg.norm(H)                      # ||H_t||_F
        lam = ratio * fro                               # adaptive ridge strength
        x = chebyshev_solve(H, q[t], lam, fro, n_iter)  # ~ (H_t + lam I)^{-1} q_t
        qm = nu[t] * x + (1 - nu[t]) * q[t]             # mix with the plain query
        out.append(U @ qm / D ** 0.5)                   # GLA-style read-out
    return torch.stack(out)

def chunk_Hx(X, K, G, H0):
    """H_t x_t for every t of one chunk without forming H_t.
    X, K: [C,D]; G: [C] chunk-local cumulative log-decay; H0: state at chunk start."""
    decay = torch.exp(G[:, None] - G[None, :]).tril()  # exp(G_t - G_j) for j <= t
    return torch.exp(G)[:, None] * (X @ H0) + ((X @ K.T) * decay) @ K
\end{lstlisting}

\subsection{Reported results}

\begin{results}
The repository reports serving-efficiency numbers for \emph{primed hybrid} models in which half of
the attention layers of Qwen3 were replaced by GKA (all from
\repofile{repos/awslabs__hybrid-model-factory/docs/Inference.md}; 8$\times$H200, tensor parallelism
8, sustained decode with a saturated KV cache):
\begin{itemize}
  \item \textbf{Decode throughput at 128K context.} GKA-primed-HQwen3-8B: 2,736 tokens/s at the
  default \code{num_iter=30}, 2.23$\times$ the Qwen3-8B baseline (1,227); 2,801 (2.28$\times$) at
  \code{num_iter=10}. For comparison the GDN-primed model reaches 2,863 (2.33$\times$). At 16K the
  GKA model gives 15,892 tokens/s (1.78$\times$) at 30 iterations and 17,261 (1.93$\times$) at 10.
  GKA-primed-HQwen3-32B reaches 1,168 tokens/s at 128K, 1.99$\times$ the Qwen3-32B baseline (586).
  \item \textbf{Time to first token} (same batch size as the saturated baseline): at 128K the 8B GKA
  hybrid needs 0.85$\times$ the Transformer's TTFT at 30 iterations and 0.82$\times$ at 10 (GDN:
  0.74$\times$); at 16K it is slower than the Transformer (1.26$\times$ at 30 iterations).
  \item \textbf{Compute control.} The document notes that most of the gain from lowering
  \code{num_iter} comes from 30 $\rightarrow$ 10 and that the effect is larger at shorter contexts,
  where the GKA layers are a larger share of the forward pass.
\end{itemize}
These numbers measure the hybrids' speed, not GKA's modelling quality. Qualitatively, the README
states that GKA is more expressive than Mamba-2 and GDN, ``translates to strong performance on
long-context and reasoning tasks'', and uniquely among the toolkit's SSM layers ``supports variable
compute at inference via its num\_iter parameter''. See \citet{peng2026gka} for accuracy results.
\end{results}

\subsection{Discussion}

\textbf{Strengths.} GKA's state is optimal for an explicit objective over the whole gated past, not
the result of a single local correction. Whitening the query cancels interference between similar
keys (\cref{ex:c08-whiten}), which is what the delta rule tries to do one token at a time. The
adaptive regularizer gives a provable bound on the condition number, which makes a fixed-interval
Chebyshev solver both sufficient and stable, and makes the iteration count a deployment-time
dial. The implementation reuses GLA kernels for everything except the solve.

\textbf{Limitations.} Training and decoding cost grow linearly with the iteration count, and the
state is twice GLA's. The solve is approximate and the regularizer always shrinks: near-duplicate
keys are separated only partially (\cref{ex:c08-whiten}). The memory is still a $d_k\times d_k$
summary, so at most $d_k$ linearly independent keys can be recalled exactly. The layer's output is
a learned mixture of the ridge read and the plain GLA read, so the regression interpretation holds
only to the extent $\nu_t\approx1$.

\textbf{Relations.} With $\mH_t+\lambda_t\mI$ replaced by $\mI$ GKA is GLA (\cref{ch:linear}); the
delta-rule family of \cref{ch:delta}, including Kaczmarz linear attention \citep{zou2026kaczmarz},
takes single projection or gradient steps toward the same regression target; and the mesa layer of
\citet{c08-vonoswald2023mesa} solves an in-context ridge regression with a Sherman--Morrison
recursion, without GKA's fading memory, adaptive regularization or iterative solver. Test-time
training layers (\cref{ch:ttt}) replace the linear map by a nonlinear one and the exact solve by
a few gradient steps.

% ===========================================================================================
\section{Next-Latent Prediction}\label{sec:c08-nextlat}

\subsection{Motivation}

KLA and GKA change \emph{how} a state is updated. Next-Latent Prediction changes \emph{what a state
should represent}, and does so without touching the architecture. A Transformer's memory is its KV
cache, which grows with the sequence, and attention can look up any past token whenever it needs
to. As the authors put it in their abstract, Transformers therefore ``lack an inherent incentive to
compress history into compact latent states with consistent transition rules'', which can lead to
solutions that generalize poorly. Recurrent models are forced to compress, but nothing forces the
compressed state to be \emph{sufficient}. \Citet{teoh2026nextlat} add one auxiliary objective to
next-token training: from the hidden state at position $t$ and the next token, predict the hidden
state at position $t+1$. A Transformer that satisfies this objective has hidden states that evolve
by a fixed transition rule and, together with the usual next-token head, form a world model whose
states are belief states (\cref{sec:c08-belief}). Architecture, parallel training and inference are
unchanged; the memory mechanism is shaped offline, by the loss.

\begin{taxonomy}
In \citet{zhoubian2026memory}, Next-Latent Prediction is \textbf{implicit} memory (hidden activations),
with \textbf{offline} update dynamics and \textbf{long-term} persistence (survey Table~I: ``Offline,
Long-Term, Latent Belief-State Induction''). Its \textbf{update rule} is the survey's
\emph{objective-induced or structural} category (Table~II), alongside Engram, MoE models and
Priming: memory behaviour is shaped by an auxiliary objective rather than by per-instance online
writes. The survey stresses that NextLat ``is therefore an offline, objective-induced mechanism
rather than a new test-time write rule'', occupying ``a different position'' from the
state-transition updates of the rest of this chapter.
\end{taxonomy}

\subsection{Formulation}\label{sec:c08-nextlat-form}

Let $x_1,\dots,x_T$ be tokens (integers in $\mathcal{V}$), $\ve(x)\in\R^d$ the token embedding,
$\vh_t\in\R^d$ the Transformer's final-layer hidden state at position $t$ (after the final norm),
and $p_{\vtheta}(x_{t+1}\mid\vh_t)=\softmax(\mW\vh_t)$ the next-token head. NextLat adds a small
\emph{latent dynamics model}
\begin{equation}\label{eq:c08-nextlat-f}
  \hat\vh_{t+1}=f_\psi\bigl(\vh_t,\ve(x_{t+1})\bigr)=\vh_t+\mathrm{MLP}_\psi\bigl(\LayerNorm([\ve(x_{t+1});\vh_t])\bigr),
\end{equation}
a residual three-layer GELU MLP on the concatenation of the next token's embedding and the current
hidden state. Writing $\mathrm{sg}(\cdot)$ for stop-gradient, the training loss is
\begin{align}
  \loss_{\mathrm{NTP}} &= -\textstyle\sum_t\log p_{\vtheta}(x_{t+1}\mid\vh_t), \nonumber\\
  \loss_{\mathrm{lat}}^{(j)} &= \textstyle\sum_t \mathrm{SmoothL1}\bigl(\hat\vh_{t+j},\,\mathrm{sg}(\vh_{t+j})\bigr), \label{eq:c08-nextlat-lat}\\
  \loss_{\mathrm{KL}}^{(j)} &= \textstyle\sum_t \mathrm{KL}\bigl(\softmax(\mathrm{sg}(\mW\vh_{t+j}))\,\|\,\softmax(\mathrm{sg}(\mW)\hat\vh_{t+j})\bigr), \label{eq:c08-nextlat-kl}\\
  \loss &= \loss_{\mathrm{NTP}}+\frac{\lambda_{\mathrm{MSE}}}{d_h}\sum_{j=1}^{d_h}\loss_{\mathrm{lat}}^{(j)}
          +\frac{\lambda_{\mathrm{KL}}}{d_h}\sum_{j=1}^{d_h}\loss_{\mathrm{KL}}^{(j)} , \label{eq:c08-nextlat-total}
\end{align}
where the multi-step predictions roll the model's own transition forward: starting from
$\hat\vh_t=\vh_t$, we set $\hat\vh_{t+j}=f_\psi(\hat\vh_{t+j-1},\ve(x_{t+j}))$ for $j=1,\dots,d_h$,
up to a horizon $d_h$ (the code's \code{mtp_horizon}). The code refers to \eqref{eq:c08-nextlat-lat} as the
hidden-state MSE term, weighted by \code{lambda_mse} (``lambda\_MSE in the paper'', says the config
comment), but computes it with PyTorch's smooth-L1 (Huber) loss. Sums are
taken as means over valid positions in the code. Each design choice has a reason.
\begin{itemize}
  \item \textbf{Stop-gradient on the target.} Gradients flow into the trunk through the
  \emph{input} $\vh_t$ (and the embedding), pushing $\vh_t$ to be predictive of $\vh_{t+1}$, but not
  through the target $\vh_{t+1}$, which would otherwise be pulled toward whatever is easy to predict.
  Collapse to a trivial constant state is prevented by $\loss_{\mathrm{NTP}}$, which still requires
  $\vh_t$ to predict the next token.
  \item \textbf{The KL term} asks that the predicted latent produce the same next-token distribution
  as the true one, i.e.\ agree with it where it matters for output, even if not in every coordinate.
  The head $\mW$ is detached on the student side, so only $\loss_{\mathrm{NTP}}$ trains the head.
  \item \textbf{Multi-step rollouts} ($d_h>1$) penalize compounding errors of the transition, as in
  model-based reinforcement learning.
\end{itemize}

\begin{figure}[t]
\centering
\begin{tikzpicture}[
  box/.style={draw,rounded corners=2pt,align=center,font=\small,minimum height=7mm,inner sep=3pt},
  arr/.style={-{Stealth[length=2mm]},thick}, darr/.style={-{Stealth[length=2mm]},thick,dashed}]
\node[box,minimum width=58mm,fill=gray!8] (trunk) {Transformer trunk (unchanged)};
\node[font=\small,below=4mm of trunk.south west,anchor=north west] (xt) {$x_{\le t}$};
\node[font=\small,below=4mm of trunk.south east,anchor=north east] (xt1) {$x_{\le t+1}$};
\draw[arr] (xt) -- (xt|-trunk.south);
\draw[arr] (xt1) -- (xt1|-trunk.south);
\node[box,above=8mm of trunk.north west,anchor=south west] (ht) {$\vh_t$};
\node[box,above=8mm of trunk.north east,anchor=south east] (ht1) {$\vh_{t+1}$};
\draw[arr] (ht|-trunk.north) -- (ht);
\draw[arr] (ht1|-trunk.north) -- (ht1);
\node[box,fill=blue!6] (f) at ($(ht)!0.5!(ht1)+(0,14mm)$) {$f_\psi$};
\node[font=\small,above=5mm of f] (e) {$\ve(x_{t+1})$};
\draw[arr] (e) -- (f);
\draw[arr] (ht) -- (f);
\node[box] (hhat) at (ht1|-f) {$\hat\vh_{t+1}$};
\draw[arr] (f) -- (hhat);
\draw[darr] (ht1) -- node[right,font=\scriptsize,align=left] {smooth-L1 to sg$(\vh_{t+1})$\\KL of the heads} (hhat);
\node[box,left=6mm of ht] (ntp) {$p(x_{t+1}\mid\vh_t)$};
\draw[arr] (ht) -- (ntp);
\end{tikzpicture}
\caption{Next-Latent Prediction. The trunk and the next-token head are trained as usual; a small
latent dynamics model $f_\psi$ predicts the next hidden state from the current one and the next
token, and its prediction is matched to the (stop-gradient) true next hidden state and to its
next-token distribution.}
\label{fig:c08-nextlat}
\end{figure}

\subsection{Why the objective produces belief states}

The following statement captures the idea; it is our informal version, and the paper proves a
precise convergence result (see the paper for its exact statement and assumptions).
\begin{proposition}[Exact next-latent prediction gives a belief state]\label{prop:c08-belief}
Suppose that for every history (a) the head is exact,
$p_{\vtheta}(x_{t+1}\mid\vh_t)=P(x_{t+1}\mid x_{1:t})$, and (b) the transition is exact,
$\vh_{t+1}=f(\vh_t,x_{t+1})$. Then for every $k\ge1$ the conditional law
$P(x_{t+1:t+k}\mid x_{1:t})$ depends on the history only through $\vh_t$: $\vh_t$ is a belief state,
and $(f,p_{\vtheta})$ is a world model in the sense of \cref{sec:c08-belief}.
\end{proposition}
\begin{proof}
By (b) and induction on $j$, $\vh_{t+j}=F_j(\vh_t,x_{t+1:t+j})$ where $F_j$ applies $f$ $j$ times.
By the chain rule and (a),
\[
  P(x_{t+1:t+k}\mid x_{1:t})=\prod_{j=1}^{k}P(x_{t+j}\mid x_{1:t+j-1})
  =\prod_{j=1}^{k}p_{\vtheta}\bigl(x_{t+j}\mid F_{j-1}(\vh_t,x_{t+1:t+j-1})\bigr),
\]
and the right-hand side is a function of $\vh_t$ and the future tokens only.
\end{proof}
Next-token prediction alone delivers (a) but not (b). A Transformer can let $\vh_t$ encode only what
is needed for $x_{t+1}$ and recover everything else later by attending to the KV cache. Condition (b)
forbids this: $\vh_{t+1}$ may contain nothing that cannot be computed from $\vh_t$ and $x_{t+1}$, so
$\vh_t$ must already carry every piece of the history the future will need. Our HMM check in
\cref{sec:c08-belief} is the same computation with the exact Bayes filter in place of $(f,\vh)$.

How does this relate to other auxiliary objectives that the repository implements as baselines?
Multi-token prediction \citep{c08-gloeckle2024mtp} adds heads that predict $x_{t+2},x_{t+3},\dots$
from $\vh_t$, which forces $\vh_t$ to contain information about the future but imposes no
consistent transition between states. The Belief State Transformer \citep{c08-hu2025bst} learns
belief states with a forward and a backward encoder that predict the next and the previous token.
NextLat obtains the recurrence constraint with one small MLP and losses in latent space.

\subsection{Self-speculative decoding}\label{sec:c08-nextlat-spec}

The latent dynamics model gives the Transformer a cheap draft model of itself. After one full
forward pass produces $\vh_t$, drafts are generated by the MLP alone:
$\hat x_{t+1}\sim\softmax(\mW\vh_t)$, $\hat\vh_{t+1}=f_\psi(\vh_t,\ve(\hat x_{t+1}))$,
$\hat x_{t+2}\sim\softmax(\mW\hat\vh_{t+1})$, and so on for $\gamma$ tokens. One full forward pass
over the prefix and the drafts then gives the target distributions $p_j$, and standard speculative
sampling \citep{c08-leviathan2023speculative} accepts draft $j$ with probability
$\min(1,p_j(\hat x)/q_j(\hat x))$, stops at the first rejection and resamples that position from
$\mathrm{normalize}(\max(0,p_j-q_j))$; if every draft is accepted it samples one bonus token from the
target. The output distribution is exactly the target model's: for one position,
\[
  P(\text{out}=y)=q(y)\min\Bigl(1,\tfrac{p(y)}{q(y)}\Bigr)
  +\Bigl(1-\sum_z\min(p(z),q(z))\Bigr)\frac{p(y)-\min(p(y),q(y))}{\sum_z\bigl(p(z)-\min(p(z),q(z))\bigr)}
  =p(y),
\]
because $1-\sum_z\min(p,q)=\sum_z(p-\min(p,q))$ (both distributions sum to one); we confirmed this
numerically to $10^{-16}$. The number of tokens accepted per round varies with how well the latent
rollout tracks the true hidden states, which is why the authors call it variable-length
self-speculative decoding. The first draft uses the true $\vh_t$ and the true head, so $q_1=p_1$ and
it is always accepted; the repository's acceptance metrics therefore exclude position~1.

\subsection{Algorithms and complexity}

\begin{algorithm}[t]
\caption{NextLat training step}\label{alg:c08-nextlat-train}
\begin{algorithmic}[1]
\Require batch of token sequences $x_{1:T}$; horizon $d_h$; weights $\lambda_{\mathrm{MSE}},\lambda_{\mathrm{KL}}$
\State $\ve_{1:T}\gets\ve(x_{1:T})$; \ $\vh_{1:T}\gets\mathrm{Trunk}(\ve_{1:T})$ \Comment{one causal forward pass}
\State $\loss\gets\loss_{\mathrm{NTP}}(\mW\vh_{1:T-1},x_{2:T})$; \ $\hat\vh\gets\vh$
\For{$j=1..d_h$}
  \State $\hat\vh_{t+1}\gets f_\psi(\hat\vh_t,\ve_{t+1})$ for all valid $t$ \Comment{parallel over positions}
  \State $\loss\gets\loss+\frac{\lambda_{\mathrm{MSE}}}{d_h}\,\mathrm{SmoothL1}(\hat\vh,\mathrm{sg}(\vh))$
  \State $\loss\gets\loss+\frac{\lambda_{\mathrm{KL}}}{d_h}\,\mathrm{KL}\bigl(\softmax(\mathrm{sg}(\mW\vh))\,\|\,\softmax(\mathrm{sg}(\mW)\hat\vh)\bigr)$
\EndFor
\State backpropagate $\loss$ into trunk, embeddings, head and $f_\psi$
\end{algorithmic}
\end{algorithm}

\begin{algorithm}[t]
\caption{NextLat self-speculative decoding (one round)}\label{alg:c08-nextlat-spec}
\begin{algorithmic}[1]
\Require prefix $x_{1:t}$; draft length $\gamma$
\State $\vh\gets\mathrm{Trunk}(x_{1:t})_t$ \Comment{full model}
\For{$j=1..\gamma$}
  \State $q_j\gets\softmax(\mW\vh)$; \ $\hat x_j\sim q_j$; \ $\vh\gets f_\psi(\vh,\ve(\hat x_j))$ \Comment{MLP only}
\EndFor
\State $p_{1..\gamma+1}\gets$ next-token distributions of the full model on $x_{1:t}\hat x_{1:\gamma}$ \Comment{one pass}
\For{$j=1..\gamma$}
  \If{$u\sim U(0,1)<\min(1,p_j(\hat x_j)/q_j(\hat x_j))$} accept $\hat x_j$
  \Else\ append a sample of $\mathrm{normalize}(\max(0,p_j-q_j))$ and \Return
  \EndIf
\EndFor
\State append a sample of $p_{\gamma+1}$
\end{algorithmic}
\end{algorithm}

\Cref{alg:c08-nextlat-train,alg:c08-nextlat-spec} summarize training and decoding. Training adds,
per horizon step, one pass of the MLP over all positions and one extra vocabulary projection (the
student logits must be materialized for the KL), i.e.\ $\bigO(d_hT(d^2+d\abs{\mathcal{V}}))$ on top of
the Transformer's cost. Inference without speculation is unchanged, and so is the memory footprint:
the KV cache still grows with $T$. NextLat reorganizes what the hidden state contains; it does not
shrink the cache. A speculative round costs one forward pass for the prefix, $\gamma$ MLP calls and
one verification pass. The repository's implementation recomputes the full prefix in both passes
(``this implementation does not include KV-cache optimization'',
\repofile{repos/JaydenTeoh__NextLat/utils/speculative_sampling.py}).

\subsection{Implementation walkthrough}

\begin{implnote}[title={In the code: model and losses}]
Everything is in \repofile{repos/JaydenTeoh__NextLat/models/model_nextlat.py}; the README points
to the \code{compute_loss()} functions of the files in \repofile{repos/JaydenTeoh__NextLat/models/}
as the clearest way to compare NextLat with its baselines (\code{model_gpt.py} for next-token
prediction, \code{model_mtp_gloeckle.py}, \code{model_mtp_jtp.py}, \code{model_bst.py}).
\begin{itemize}
  \item \code{NextLatConfig}: \code{lambda_kl}, \code{lambda_mse}, \code{lambda_ce} (an optional
  cross-entropy on the next-next token, 0 by default and used for logging), \code{mtp_horizon}
  ($d_h$) and \code{proj_factor}.
  \item \code{NextLatDynamicsModel} is $f_\psi$: input width $2d$, hidden width
  $128\cdot\mathrm{round}(\text{\code{proj_factor}}\cdot2d/128)$ (for $d=768$ and factor 1.3 this is 2048),
  \code{LayerNorm} on the concatenation \texttt{[next\_token\_embeds,~hidden\_states]}, and a residual
  output (\cref{lst:c08-nextlat-f}).
  \item \texttt{NextLatTransformer.forward(...,~return\_hidden\_states=True)} returns the token
  embeddings and the normalized final hidden states, the two inputs of the latent losses.
  \item \code{_nextlat_compute_losses} implements the horizon loop: at each step it shifts the
  predicted states left and the embeddings, targets and teacher logits right by one, applies
  \code{dynamics_model}, and calls \code{_nextlat_loss_function} (\cref{lst:c08-nextlat-loss}). The
  smooth-L1 term skips targets that are end-of-sequence tokens (they start a new document in packed
  sequences), and the token-level KL skips next-next tokens that cross a document boundary.
  \item \code{compute_loss} computes all losses on \emph{detached} copies of the hidden states and
  embeddings (re-marked \code{requires_grad}), backpropagates the loss to those copies, and then
  calls \texttt{fabric.backward(combined\_emb,~gradient=combined\_grad)} to push the accumulated
  gradients through the trunk once. The comments explain the purpose: the loss function can be
  compiled separately with \code{torch.compile}, and gradients from all losses are combined before
  the single backward pass through the trunk.
\end{itemize}
\end{implnote}

\begin{lstlisting}[caption={Excerpt from \texttt{repos/\allowbreak JaydenTeoh\_\_NextLat/\allowbreak models/\allowbreak model\_nextlat.py} (\texttt{NextLatDynamicsModel.forward}).},label={lst:c08-nextlat-f}]
        hidden_states = self.hidden_state_dropout(current_states)
        x = torch.cat([next_token_embeds, hidden_states], dim=-1)
        x = self.norm_x(x)

        # residual connection
        delta = self.mlp(x)
        next_states = delta + current_states

        return next_states
\end{lstlisting}

\begin{lstlisting}[caption={Excerpt from \texttt{repos/\allowbreak JaydenTeoh\_\_NextLat/\allowbreak models/\allowbreak model\_nextlat.py} (\texttt{NextLat.\_nextlat\_loss\_function}).},label={lst:c08-nextlat-loss}]
        mse_elem = F.smooth_l1_loss(pred_h_t_next, h_t_next.detach(), reduction="none")
        w = mse_mask.unsqueeze(-1).to(dtype=mse_elem.dtype)
        # Same as reduction="mean" over masked (B, T, n_embd) elements.
        # i.e., divide over B*T*n_embd elements.
        denom = w.expand_as(mse_elem).sum().clamp_min(1.0)
        MSE = (mse_elem * w).sum() / denom
        # ...
        pred_token_inputs = pred_h_t_next[:, :-1]
        # Keep NextLat CE/KL from updating lm_head by using detached head weights.
        # Have to materialize student logits to compute KL loss
        lm_head_weight = self.model.lm_head.weight.detach()
        pred_tokens = F.linear(pred_token_inputs, lm_head_weight)
        # ...
        # Detach teacher logits to avoid updating teacher computation graph
        kl_loss = self._categorical_kl_loss(
            teacher_logits.detach(), pred_tokens, token_pred_mask
        )
\end{lstlisting}

Here \code{teacher_logits} are the next-token logits $\mW\vh$ computed for $\loss_{\mathrm{NTP}}$,
shifted so that position $t$ holds $\mW\vh_{t+1}$, and \code{_categorical_kl_loss} computes
$\mathrm{KL}(\text{teacher}\,\|\,\text{student})$ per position with \code{F.kl_div} and averages over
valid positions.

\begin{implnote}[title={In the code: drafting, verification, configs}]
\code{NextLat.speculative_propose} (\cref{lst:c08-nextlat-propose}) runs the full model once on the
prefix, then alternates head sampling and \code{dynamics_model} steps. Verification and the
accept/resample loop are in \code{speculative_decode_v2} in
\repofile{repos/JaydenTeoh__NextLat/utils/speculative_sampling.py} (adapted, as its docstring says,
from an open-source implementation of Leviathan et al.'s algorithm), with target probabilities from
\code{build_speculative_target_probs_fn} in \repofile{repos/JaydenTeoh__NextLat/models/model_speculative.py}.
Benchmarks are driven by \code{eval/eval_speculative_checkpoints.py} (\code{--gamma} sets the draft
length). Configs under \repofile{repos/JaydenTeoh__NextLat/config/} cover Manhattan taxi rides,
Countdown, path-star graphs, TinyStories and FineWeb-Edu; for example
\code{config/fineweb/100M/nextlat_finewebedu.yaml} trains a 12-layer, 12-head, $d=768$ model with
context 1024, $\lambda_{\mathrm{KL}}=\lambda_{\mathrm{MSE}}=1$, horizon 1, \code{proj_factor} 1.3, the
Muon optimizer, about 10B tokens (19,074 updates of 524,288 tokens), and in-training speculative
evaluation with draft length 4. The 1B FineWeb-Edu configs come in horizon-1 and horizon-2 variants.
\end{implnote}

\begin{lstlisting}[caption={Excerpt from \texttt{repos/\allowbreak JaydenTeoh\_\_NextLat/\allowbreak models/\allowbreak model\_nextlat.py} (\texttt{NextLat.speculative\_propose}).},label={lst:c08-nextlat-propose}]
        core = self._unwrap_module()
        _, hidden_states = core(seq, return_hidden_states=True)
        state = hidden_states[:, -1, :]
        drafted: List[torch.Tensor] = []
        q_probs_steps: List[torch.Tensor] = []

        for _ in range(steps_to_propose):
            logits = core.lm_head(state)
            q_probs = normalize_logits(
                logits, temperature=temperature, top_k=top_k, top_p=top_p
            )
            tok = sample_from_probs(q_probs)
            drafted.append(tok)
            q_probs_steps.append(q_probs)
            next_token_emb = core.token_embedding(tok).squeeze(1)
            state = core.dynamics_model(state, next_token_emb)

        return torch.cat(drafted, dim=1), q_probs_steps
\end{lstlisting}

\subsection{Reference implementation}

\Cref{lst:c08-nextlat-ours} gives the loss \eqref{eq:c08-nextlat-total} and the drafting loop in a
few lines, for a generic causal trunk. The index bookkeeping follows the repository's: at horizon
step $j$ (0-based) the predicted states are $\hat\vh_{2+j},\dots,\hat\vh_T$ and the teacher logits are
$\mW\vh_{2+j},\dots,\mW\vh_{T-1}$. We could not execute this listing (PyTorch is not installed in
our environment); its probabilistic content, the belief-state rollout and the exactness of the
acceptance rule, was checked numerically as described above.

\begin{lstlisting}[caption={Reference implementation (ours): the NextLat loss and latent drafting.},label={lst:c08-nextlat-ours}]
import torch, torch.nn as nn, torch.nn.functional as F

class LatentDynamics(nn.Module):              # f_psi(h_t, e(x_{t+1})) -> h_{t+1}
    def __init__(self, d, proj=1.3):
        super().__init__()
        hid = 128 * round(proj * 2 * d / 128)
        self.norm = nn.LayerNorm(2 * d)
        self.mlp = nn.Sequential(nn.Linear(2 * d, hid), nn.GELU(),
                                 nn.Linear(hid, hid), nn.GELU(), nn.Linear(hid, d))
    def forward(self, h, e_next):
        return h + self.mlp(self.norm(torch.cat([e_next, h], -1)))   # residual

def nextlat_loss(trunk, emb, head, f, x, horizon=1, l_mse=1.0, l_kl=1.0):
    """x: [B,T] tokens; trunk maps embeddings to final (normed) hidden states."""
    e = emb(x); h = trunk(e)                          # [B,T,d]
    logits = head(h[:, :-1])                          # predicts x_2..x_T
    loss = F.cross_entropy(logits.flatten(0, 1), x[:, 1:].flatten())
    teacher = F.log_softmax(logits.detach(), -1)      # log p(. | h_t), no grad
    h_pred = h
    for j in range(horizon):
        h_pred = f(h_pred[:, :-1], e[:, j + 1:])      # roll the latent forward
        loss = loss + l_mse / horizon * F.smooth_l1_loss(h_pred, h[:, j + 1:].detach())
        student = F.log_softmax(F.linear(h_pred[:, :-1], head.weight.detach()), -1)
        loss = loss + l_kl / horizon * F.kl_div(       # KL(teacher || student)
            student.flatten(0, 1), teacher[:, j + 1:].flatten(0, 1),
            log_target=True, reduction='batchmean')
    return loss

@torch.no_grad()
def draft(trunk, emb, head, f, prefix, gamma):        # cheap self-drafting
    h = trunk(emb(prefix))[:, -1]                     # one full forward pass
    toks, qs = [], []
    for _ in range(gamma):
        q = F.softmax(head(h), -1); tok = torch.multinomial(q, 1)
        toks.append(tok); qs.append(q)
        h = f(h, emb(tok).squeeze(1))                 # MLP step, no attention
    return torch.cat(toks, 1), qs
\end{lstlisting}

\subsection{Reported results}

\begin{results}
The repository contains no result tables; its README describes the evaluations without numbers
(\repofile{repos/JaydenTeoh__NextLat/README.md}):
\begin{itemize}
  \item \textbf{Manhattan taxi rides}: generated trajectories are used to reconstruct the model's
  internal map of Manhattan and to score sequence compression, detour robustness and the effective
  rank of the latents (\code{data/manhattan/}).
  \item \textbf{TinyStories}: frozen models are probed with linear probes on the hidden states to
  predict 1--20 tokens ahead.
  \item \textbf{FineWeb-Edu pretraining}: downstream benchmarks via the LM Evaluation Harness, and
  self-speculative decoding of the multi-token-prediction models (JTP, MTP, NextLat) across domains,
  reporting acceptance rate and accepted tokens per step.
  \item Path-Star graphs and Countdown accuracies are tracked during training.
\end{itemize}
The survey \citep{zhoubian2026memory} summarizes the mechanism as an objective ``encouraging hidden
representations to preserve history useful for future prediction''. See \citet{teoh2026nextlat}
for the reported gains in accuracy, representation compression, planning and decoding speed.
\end{results}

\subsection{Discussion}

\textbf{Strengths.} NextLat changes neither the architecture nor inference, so it composes with any
Transformer recipe; the extra cost is a small MLP and a few losses. It gives a concrete,
testable meaning to ``compact world model'' (\cref{prop:c08-belief}), and the same MLP doubles as a
free draft model.

\textbf{Limitations.} The guarantee holds only at the optimum of the objective; in practice the
transition is approximate and errors compound over long rollouts, which is what the horizon
$d_h$ and the KL term try to control. The KV cache, the actual memory footprint, is unchanged: a
belief-state Transformer could in principle drop old cache entries, but NextLat itself does not.
The extra vocabulary-sized projection for the KL term is not free for large vocabularies.

\textbf{Relations.} NextLat is the ``objective-induced'' counterpart of this chapter's other systems:
KLA and GKA hard-code an estimator into the state update, while NextLat asks gradient descent to
discover a state with the defining property of an estimator's state. It connects to compressive or
consolidating designs that do shrink the cache (Bottlenecked Transformers, \cref{ch:routing},
\citep{oomerjee2026bottlenecked}) and to hybrids that pair a recurrent state with attention
(\cref{ch:hybrid}).

% ===========================================================================================
\section{Comparison and discussion}\label{sec:c08-comparison}

\begin{table}[t]
\centering
\small
\caption{The three systems of this chapter side by side ($d_v\times d_k$ state per head or layer;
$n_\text{iter}$ solver iterations; $\gamma$ draft length).}
\label{tab:c08-compare}
\begin{tabularx}{\textwidth}{@{}>{\raggedright\arraybackslash}p{2.3cm}>{\raggedright\arraybackslash}X>{\raggedright\arraybackslash}X>{\raggedright\arraybackslash}X@{}}
\toprule
 & \textbf{Kalman Linear Attention} & \textbf{Gated KalmaNet} & \textbf{Next-Latent Prediction} \\
\midrule
What is estimated & Posterior over every state cell of a factorized linear-Gaussian model & Ridge-regression map $\mS_t$ over the gated past & Nothing explicitly; hidden state trained to be a belief state \\
State & $(\mLambda_t,\mU_t)$: $2d_vd_k$ & $(\mH_t,\mU_t)$: $2d_k^2$ per head & Unchanged KV cache \\
Write & Precision M\"obius update; GLA-like write gated by $a/(a^2+p\Lambda_{t-1})$ & GLA writes to $\mH_t$ (values $=$ keys) and $\mU_t$ & No new write rule \\
Read & $(\mU_t\oslash\mLambda_t)\vq_t$, plus variance & $\mU_t(\mH_t+\lambda_t\mI)^{-1}\vq_t$, mixed with GLA read & Standard attention \\
Uncertainty & Explicit per-cell variance & Implicit (posterior covariance never formed) & None explicit \\
Forgetting & Process noise $p$ caps precision & Decay gate $\alpha_t$ fades $\mH_t,\mU_t$; the ridge $\lambda_t$ is not decayed & Not applicable \\
Parallel training & Two associative scans & Chunked Chebyshev + chunkwise GLA & Standard Transformer \\
Decode cost/token & $\bigO(d_vd_k)$ & $\bigO(n_\text{iter}d_k^2)$ per head & Transformer cost; optional drafting with $\gamma$ MLP steps \\
Survey update rule & State transition (filtering) & State transition (online regression) & Objective-induced (offline) \\
Code (mirror) & \texttt{repos/\allowbreak vaisakh-shaj\_\_kalman-\allowbreak linear-attention} & \texttt{repos/\allowbreak awslabs\_\_hybrid-\allowbreak model-factory} & \texttt{repos/\allowbreak JaydenTeoh\_\_NextLat} \\
\bottomrule
\end{tabularx}
\end{table}

\Cref{tab:c08-compare} puts the three systems side by side. KLA and GKA sit at opposite corners of
the same design space. KLA keeps \emph{dynamics} (a transition and process noise) and pays for
tractability with a fully factorized posterior: it knows how uncertain each cell is but ignores
correlations between cells. GKA keeps the \emph{full key covariance}, and with it the ability to
cancel interference between correlated keys, but assumes a static map with fading data in place of
process noise, and never forms the covariance explicitly: it touches it only through linear solves.
KLA buys its parallelism with algebra (the information form turns the filter into scans); GKA buys
it with numerics (an adaptive regularizer makes a fixed-coefficient solver safe inside a chunked
kernel). NextLat stands apart: it adds no state and no write rule, and instead uses the training
objective to make an existing representation carry the property, sufficiency, that the other two
build in by construction.

Practical consequences follow. KLA is the natural choice when per-token uncertainty is itself
useful downstream, or when decoding cost must stay at GLA levels. GKA is the natural choice for
recall-heavy workloads inside hybrids, where its extra solve buys precision and its iteration count
can be tuned at deployment (\cref{ch:hybrid}). NextLat can be added to any Transformer training run,
including those of hybrid models, and yields a draft model as a by-product.

% ===========================================================================================
\section{Memory as estimation}\label{sec:c08-estimation}

The background section showed that the delta rule, RLS and the Kalman filter are three ways of
estimating the same regression map. We can now place the book's recurrent memories on a single
scale, ordered by how much of the estimation problem each one solves (\cref{tab:c08-ladder}).

\begin{table}[t]
\centering
\small
\caption{Recurrent memories as estimators of the map $\mS$ with $\mS\vk_i\approx\vv_i$. ``Exact''
means the minimizer of the stated objective (up to solver tolerance for GKA).}
\label{tab:c08-ladder}
\begin{tabularx}{\textwidth}{@{}>{\raggedright\arraybackslash}p{3.0cm}>{\raggedright\arraybackslash}X>{\raggedright\arraybackslash}p{3.1cm}>{\raggedright\arraybackslash}p{2.7cm}@{}}
\toprule
\textbf{Memory} & \textbf{Objective and step} & \textbf{Preconditioner} & \textbf{Cost/token} \\
\midrule
Linear attention / GLA (\cref{ch:linear}) & One unit step on $-\inner{\mS\vk_t}{\vv_t}$ (plus decay) & none & $\bigO(d_kd_v)$ \\
Delta rule, Gated DeltaNet (\cref{ch:delta}) & One SGD step on $\tfrac12\norm{\mS\vk_t-\vv_t}^2$ from $\mS_{t-1}$ & $\beta_t\mI$ & $\bigO(d_kd_v)$ \\
Kaczmarz LA (\cref{ch:delta}) & Projection onto $\{\mS:\mS\vk_t=\vv_t\}$ & $\mI/\norm{\vk_t}^2$ & $\bigO(d_kd_v)$ \\
RLS / static Kalman filter & Exact minimizer of $\sum_{i\le t}\norm{\mS\vk_i-\vv_i}^2+\lambda\norm{\mS}^2$ (one Newton step) & $\mP_t$, i.e.\ $(\mH_t+\lambda\mI)^{-1}$ & $\bigO(d_k^2+d_kd_v)$ \\
Gated KalmaNet & Exact minimizer of the gated, adaptively regularized objective \eqref{eq:c08-gka-obj}, read along $\vq_t$ only & $(\mH_t+\lambda_t\mI)^{-1}$ via Chebyshev & $\bigO(n_\text{iter}d_k^2)$ \\
Kalman Linear Attention & Exact posterior of a factorized dynamic model with process noise & per-cell $1/\Lambda_t$ & $\bigO(d_kd_v)$ \\
Test-time training (\cref{ch:ttt}) & A few gradient steps on a nonlinear memory's loss & optimizer-dependent & model-dependent \\
Next-Latent Prediction & Changes what is estimated: trains $\vh_t$ to be a sufficient statistic & --- & unchanged \\
\bottomrule
\end{tabularx}
\end{table}

Three observations organize the table.

\paragraph{The preconditioner is a per-direction learning rate.} Every row of the form
$\mS_t=\mS_{t-1}+(\vv_t-\mS_{t-1}\vk_t)\vk_t\T\bm{B}_t$ differs only in the matrix $\bm{B}_t$. The
delta rule uses a scalar. Kaczmarz normalizes by the key norm, which makes each write an exact
projection. RLS uses the inverse key covariance, which gives each direction a learning rate
inversely proportional to how much evidence it has already received; this is the
uncertainty-driven step size of the scalar example in \cref{sec:c08-kf}. Per-channel gates in
Kimi Delta Attention or Gated DeltaNet-2 (\cref{ch:delta}) can be read as cheap diagonal
approximations of this preconditioner. KLA's per-cell gate $\alpha_t$ and normalizer $1/\Lambda_t$
are another diagonal choice, derived rather than learned.

\paragraph{Exactness costs compute or structure.} Solving the regression exactly at every token
requires either a covariance-sized state that is updated with care (RLS, which is fragile with
forgetting and in low precision), an iterative solve per token (GKA, $n_\text{iter}$ matrix--vector
products), or a factorization that makes the problem small (KLA). The first-order rules are cheap
precisely because they solve only the newest term of the objective. Which end of the scale is better
is an empirical question: keys, values and gates are learned, so the network can compensate for a
crude estimator, and an exact solution of the wrong objective is not obviously better than an
approximate solution of the right one.

\paragraph{From updating a state to shaping one.} All rows but the last fix a state and an update
rule and differ in how faithfully the rule estimates. NextLat inverts the question: it keeps the
Transformer's machinery and asks training to produce a representation with the defining property of
an estimator's state. The two approaches are complementary; nothing prevents training a KLA or GKA
hybrid with a next-latent objective.

\begin{summarybox}
\begin{itemize}
  \item A key--value memory can be viewed as an estimator of the map $\mS$ with $\mS\vk_i\approx\vv_i$.
  The delta rule is one SGD step on the newest pair; RLS, equivalently the posterior mean of a static
  Kalman filter, is the exact ridge solution, obtained by one Newton step whose preconditioner is the
  posterior covariance. The delta rule is RLS with that covariance replaced by $\beta_t\mI$.
  \item In the information form of the Kalman filter, measurement updates add, and prediction is
  linear (no process noise) or linear-fractional (scalar state with process noise). Both kinds of maps
  compose associatively, so filtering becomes a parallel scan.
  \item \textbf{Kalman Linear Attention} runs an exact, factorized Kalman filter in every state cell:
  a M\"obius scan for the precisions and an affine (GLA-like) scan for the information. Its write
  gate $a/(a^2+p\Lambda_{t-1})$ comes from the accumulated precision, its read-out divides by the
  precision, it returns a variance for free, and process noise is its forgetting mechanism.
  \item \textbf{Gated KalmaNet} sets the state to the solution of a gated, fading ridge regression and
  reads it as $\mU_t(\mH_t+\lambda_t\mI)^{-1}\vq_t$: a whitened query read from a GLA memory. The
  adaptive ridge $\lambda_t=r\norm{\mH_t}_F$ bounds the condition number by $1+1/r$, which makes a
  fixed-interval Chebyshev solver stable in low precision and turns the iteration count into a
  test-time compute dial; chunked kernels evaluate $\mH_t\vx$ without forming $\mH_t$.
  \item \textbf{Next-Latent Prediction} adds a latent dynamics MLP and losses that predict the next
  hidden state; at the optimum, hidden states are belief states. It is offline, objective-induced
  memory: no architectural change, an unchanged KV cache, and a free draft model for exact
  self-speculative decoding.
\end{itemize}
\end{summarybox}

\section*{Exercises}

\begin{exercise}[Information-form filter for one cell]\label{ex:c08-ex-info}
(a) Starting from the scalar model $z_t=az_{t-1}+w_t$, $y_t=cz_t+e_t$ with $\mathrm{Var}(w_t)=p$ and
$\mathrm{Var}(e_t)=\sigma^2$, derive \eqref{eq:c08-kla-prec} and \eqref{eq:c08-kla-info} with
$\phi_t=c^2/\sigma^2$ and $r_t=cy_t/\sigma^2$. (b) Show that the derivative of the map
$\Lambda\mapsto\Lambda/(a^2+p\Lambda)+\phi$ equals $\alpha^2$ with $\alpha=a/(a^2+p\Lambda)$, and
conclude that for constant $\phi>0$ and $p>0$ the precision converges to the fixed point
$\Lambda^\star$. (c) Show that rescaling the leaf matrix \eqref{eq:c08-kla-leaf} by any $c>0$ leaves
$\Lambda_t$ unchanged, and that the product of two leaves has positive trace.
\end{exercise}

\begin{exercise}[RLS, ridge and the delta rule]\label{ex:c08-ex-rls}
(a) Prove the Sherman--Morrison formula $(\bm{M}+\vk\vk\T)^{-1}=\bm{M}^{-1}-\bm{M}^{-1}\vk\vk\T
\bm{M}^{-1}/(1+\vk\T\bm{M}^{-1}\vk)$ and use it to show by induction that RLS with
$\mP_0=\lambda^{-1}\mI$, $\mS_0=\bm{0}$ produces the ridge solution \eqref{eq:c08-ridge} at every $t$.
(b) Show that RLS with forgetting factor $\gamma$ (replace $\mP_{t-1}$ by $\mP_{t-1}/\gamma$) minimizes
$\sum_i\gamma^{t-i}\norm{\mS\vk_i-\vv_i}^2+\lambda\gamma^t\norm{\mS}_F^2$, and explain why the vanishing
regularizer is a problem when keys stop exciting some direction. (c) Write the delta rule as RLS with
$\mP_t$ replaced by $\beta_t\mI$ and find the $\beta_t$ for which one delta step from $\mS_{t-1}$ makes
$\mS_t\vk_t=\vv_t$ exactly (this is the Kaczmarz step of \cref{ch:delta}).
\end{exercise}

\begin{exercise}[How many Chebyshev iterations?]\label{ex:c08-ex-cheb}
(a) Using \cref{prop:c08-kappa} and \eqref{eq:c08-cheb-poly}, compute the number of iterations needed
to guarantee relative error $10^{-3}$ for $r\in\{0.005,0.02,0.1\}$, and the number gradient descent
with the optimal step would need. (b) Redo \cref{ex:c08-whiten} for these values of $r$ and plot the
retrieved value against $r$. What does the regularizer trade against what? (c) Explain why the
Chebyshev output after $n$ steps is a linear function of $\vq$ with a symmetric Jacobian, and why
this lets the backward pass reuse the forward solver, as in \eqref{eq:c08-gka-bwd}.
\end{exercise}

\begin{exercise}[Coding: chunked GKA]\label{ex:c08-ex-code}
Implement \cref{alg:c08-gka-prefill} in NumPy for one head, using \eqref{eq:c08-gka-chunkHX} and
\eqref{eq:c08-gka-fro}, and check it against \code{gka_recurrent} of \cref{lst:c08-gka-ours} for
several chunk sizes. Then implement the backward pass for $\tilde\vq_t$ by \eqref{eq:c08-gka-bwd}
(including the term from $\lambda_t=r\norm{\mH_t}_F$) and verify it against finite differences.
\end{exercise}

\begin{exercise}[Belief states need the transition constraint]\label{ex:c08-ex-belief}
(a) Extend the proof of \cref{prop:c08-belief} to show that if (a) and (b) hold then any two histories
with the same $\vh_t$ have identical conditional distributions over all futures. (b) Construct a
two-symbol source and two histories for which $P(x_{t+1}\mid\text{history})$ is identical but
$P(x_{t+2}\mid\text{history},x_{t+1})$ differs. Explain why a Transformer trained only on next-token
prediction may map both histories to the same $\vh_t$, and why condition (b) forbids it.
\end{exercise}

\begin{exercise}[Speculative sampling]\label{ex:c08-ex-spec}
(a) Prove that the accept/resample rule of \cref{sec:c08-nextlat-spec} outputs exactly $p$, and that
the probability of accepting a draft is $\sum_y\min(p(y),q(y))=1-\mathrm{TV}(p,q)$. (b) If drafts at
positions $2,\dots,\gamma$ are accepted independently with probability $\pi$, compute the expected
number of tokens produced per verification pass. Why is the first draft of NextLat always accepted?
\end{exercise}
